% Generated by code/build.py from the modular article.

% BEGIN preamble.tex
\documentclass[11pt]{article}
\usepackage[T1]{fontenc}
\usepackage{lmodern}
\usepackage[margin=1in]{geometry}
\usepackage{amsmath,amssymb,amsthm,mathtools,mathrsfs}
\usepackage{graphicx,booktabs,longtable,array,enumitem,xcolor,microtype,xurl}
\usepackage[colorlinks=true,linkcolor=blue!45!black,citecolor=blue!45!black,urlcolor=blue!55!black]{hyperref}
\usepackage{bookmark,fancyhdr}
\emergencystretch=3em
\setcounter{tocdepth}{2}
\numberwithin{equation}{section}
\theoremstyle{plain}
\newtheorem{theorem}{Theorem}[section]
\newtheorem{proposition}[theorem]{Proposition}
\newtheorem{lemma}[theorem]{Lemma}
\newtheorem{corollary}[theorem]{Corollary}
\newtheorem{conjecture}[theorem]{Conjecture}
\theoremstyle{definition}
\newtheorem{definition}[theorem]{Definition}
\newtheorem{example}[theorem]{Example}
\theoremstyle{remark}
\newtheorem{remark}[theorem]{Remark}
\newtheorem{question}{Question}
\DeclareMathOperator{\Li}{Li}
\DeclareMathOperator{\Log}{Log}
\DeclareMathOperator{\Res}{Res}
\DeclareMathOperator{\CT}{CT}
\DeclareMathOperator{\FP}{FP}
\DeclareMathOperator{\sgn}{sgn}
\DeclareMathOperator{\lcm}{lcm}
\newcommand{\C}{\mathbb C}
\newcommand{\R}{\mathbb R}
\newcommand{\N}{\mathbb N}
\newcommand{\Z}{\mathbb Z}
\newcommand{\Q}{\mathbb Q}
\newcommand{\ii}{\mathrm i}
\newcommand{\dd}{\,\mathrm d}
\newcommand{\eps}{\varepsilon}
\newcommand{\src}[1]{\nolinkurl{#1}}
\newcommand{\shuffle}{\mathbin{\amalg}}
\newcommand{\ordder}[1]{^{[#1]}}
\newcommand{\Hp}{\mathsf H}
\newcommand{\Mcal}{\mathcal M}
\newcommand{\Tcal}{\mathcal T}
\pagestyle{fancy}
\fancyhf{}
\fancyhead[L]{\small Independent orders and exact reductions}
\fancyhead[R]{\small ProveIt research continuation}
\fancyfoot[C]{\thepage}
\setlength{\headheight}{14pt}
\newcolumntype{L}[1]{>{\raggedright\arraybackslash}p{#1}}
\hypersetup{pdftitle={Independent Orders and Exact Reductions in Polylogarithm--Zeta Calculus},pdfauthor={ProveIt research continuation}}

% END preamble.tex

\title{Independent Orders and Exact Reductions\\in Polylogarithm--Zeta Calculus\\[7pt]
\large Entire Hurwitz interval functions, cubic Tornheim directions,\\mixed inner orders, and rigidity of Fourier identities}
\author{Research continuation prepared for Vladimir Reshetnikov\\\small ProveIt project}
\date{11 October 2026 (UTC)}
\begin{document}
\maketitle

% BEGIN sections/00_abstract.tex
\begin{abstract}
We continue the exact-identity program in the ProveIt manuscript and its
incoming research reports at revision
\texttt{5a790187c8e186e41e2b990b4941cb7a1a3c7b6b}.
Four related questions are treated with complete proofs.
First, the independent-order harmonic normalization is identified with
the established interpolant of Ihara, Nakamura, and Yamamoto. An explicit
prefix-residue proof gives its two-endpoint version, joint entire
continuation in all orders, composition and shift laws, Gamma generating
functions, Stieltjes order jets, and normalized primitives.
Second, every admissible third derivative of
$T(at,bt,ct)$ at zero is reduced to ordinary zeta derivatives and two
specified scalar coordinates: the diagonal derivative $\Omega$ and an
absolutely convergent logarithmic Gamma series $\kappa$. Exact conic and
cyclic combinations eliminate the latter; the full boundary germ is
also evaluated by a convergent Gamma series.
Third, arbitrary weighted colored sums with any mixture of positive
and nonpositive integral inner orders admit a terminating finite
reduction to nested polylogarithms with one free complex outer order.
The construction handles coincident colors, preserves the cumulative
alphabet, controls depth, and gives all Laurent constants at
nonpositive outer orders by finite local coefficient extraction.
Finally, a sampling uniqueness theorem proves the exact finite
single-centre reduction criterion for any number of shifted powers,
including identities on integer Fourier modes.
We distinguish repository advances from known theorems and from
unevaluated coordinates. The Gaussian $S_6$ and revised $S_8$ candidates
remain conjectural. Exact algebra, independent numerical diagnostics,
an executable mixed-order compiler, integration notes, and specific
further research questions accompany the article.
\end{abstract}

\noindent\textbf{Keywords.}
Multiple polylogarithms; Mordell--Tornheim zeta functions; independent
spectral orders; truncated multiple zeta interpolation; generalized
Stieltjes constants; Gamma functions; harmonic sums; Laurent constants;
Fourier uniqueness.

% END sections/00_abstract.tex

\clearpage
\tableofcontents
\clearpage

% BEGIN sections/01_scope.tex
\section{Research scope, status, and conventions}
\label{scope:section}

\subsection{The source state and the questions answered}

The source snapshot is the public ProveIt repository at commit
\begin{center}
\texttt{5a790187c8e186e41e2b990b4941cb7a1a3c7b6b},
\end{center}
inspected on 11 October 2026 (UTC). The two requested directories are
\src{Analysis/Polylogarithms/docs/manuscript} and \src{docs/incoming}
\cite{RepoMain,RepoIncoming}. All eleven incoming archives in that
snapshot were inventoried and extracted. The mathematical review
concentrated on the relevant harmonic, directional-zeta, mixed-spectral,
and Gaussian-certificate sections; this is not a claim to have
independently verified every page of the multi-volume manuscript.

The latest incoming reports already contain substantial results.
The diagonal cubic digamma finite part is reduced to a third diagonal
Tornheim derivative, all quadratic Tornheim directions are known, the
all-nonpositive inner-order sector has a finite rational reduction,
and common-shift harmonic Laurent formulas and Newton interpolations
are available. Repeating those theorems would not address the next
questions. Table~\ref{scope:map} records the precise continuations here.

\begin{table}[htbp]
\centering\small
\begin{tabular}{@{}L{0.28\textwidth}L{0.43\textwidth}L{0.22\textwidth}@{}}
\toprule
Source question & Answer in this article & Status \\
\midrule
Independent harmonic exponents (Exact Identities, Question 8)
& Entire ordered normalization, explicit relative endpoints and
  spectral consequences; Section~\ref{eht:sec}
& Analytic-normalization component answered; full question remains open \\
\addlinespace
All asymmetric cubic Tornheim rays (directional report, Question 2)
& Formula for every admissible direction in two specified scalar
  coordinates; Section~\ref{tr3:section}
& Proved reduction; no independence or ordinary-constant closure claimed \\
\addlinespace
Mixed inner signs at arbitrary depth (directional report, Question 11)
& Finite algorithm retaining a free outer order and repeated colors;
  Section~\ref{mixed:section}
& Proved; exact compiler supplied \\
\addlinespace
Scalar identities from integer Fourier samples (mixed-spectral report)
& Uniqueness on logarithmically dense samples and complete multicentre
  criterion; Section~\ref{sampling:section}
& Proved in the stated finite power--logarithm class \\
\addlinespace
Short Gaussian $S_6$ and revised $S_8$
& Search scope and unchanged conjectural status;
  Section~\ref{audit:section}
& Unresolved \\
\bottomrule
\end{tabular}
\caption{Contribution map. ``Proved'' refers to the stated theorem,
not to a claim of literature-wide priority.}\label{scope:map}
\end{table}

The literature reconciliation is substantive. The one-endpoint
entire interpolant is precisely the function constructed by Ihara,
Nakamura, and Yamamoto \cite{eht:INY}; agreement is established by
their finite reversed-star formula, not by guessing from integer
values. Its identification supplies the independent-order entire
normalization requested in one component of Question 8 of the
exact-identities report \cite{RepoExact}. That question also asks for
a smallest generating kernel with a Gamma-separated singular part;
neither minimality nor that additional factorization is established
here. The additional calculations
are presented as exact consequences and useful research organization.
The mixed-order and Tornheim theorems are advances relative to the
specified source snapshot. The literature search was targeted and
does not establish an absolute priority claim for every formula.

\subsection{What a reduction means}

A reduction is stated relative to a specified output class. A finite
sum of multiple polylogarithms with a free complex first index is an
exact and useful answer even when it does not collapse to ordinary
Gamma and zeta values. Likewise, expressing all cubic directions in
$\Omega$ and $\kappa$ is a theorem about a finite spanning collection;
it is not a theorem that these coordinates are irreducible or
arithmetically independent. In the harmonic sector, the separate
order jets retain the nonsymmetric coordinate $\eta(a)$ from the
ordered-resonance report \cite{RepoOrdered}.

These distinctions prevent three common overstatements: treating a
new name for an integral as its evaluation, treating a finite
normal form as a minimal basis, and treating equality at special
integer parameters as permission to differentiate in a transverse
parameter. Every differentiation below is performed on an identity
valid in a neighborhood in the variable being differentiated.

\subsection{Conventions needed throughout}

The superscript $\zeta^{(j)}(s,a)$ denotes differentiation in the
first, spectral variable. The generalized Stieltjes constants are
fixed by
\begin{equation}\label{scope:Stieltjes}
 \zeta(1+u,a)=\frac1u+
       \sum_{j\ge0}\frac{(-1)^j\gamma_j(a)}{j!}u^j,
 \qquad \gamma_0(a)=-\psi(a).
\end{equation}
Write $\gamma_j=\gamma_j(1)$ and $\gamma=\gamma_0$.
For a finite upper bound, the strict harmonic convention is
\[
 H_N(s_1,\ldots,s_r)
   =\sum_{N\ge n_1>\cdots>n_r\ge1}
                          n_1^{-s_1}\cdots n_r^{-s_r},
 \qquad H_N^{(s)}=H_N(s).
\]
The empty word has value one. The local Hurwitz notation in
Section~\ref{eht:sec} starts its shifted summation at zero, and its
finite-interval formula specifies the conversion explicitly.
Bernoulli numbers use $B_1=-1/2$ and
$te^{xt}/(e^t-1)=\sum_{n\ge0}B_n(x)t^n/n!$.

The notation $\CT_{C=c}$ always means the Laurent constant in the
displayed affine coordinate $C-c$. It does not mean substitution of
individual singular summands. A derivative of a restricted Tornheim
ray means that the ray is formed first and its resulting one-variable
meromorphic germ is then continued and differentiated. Coherent
principal branches are used initially in right half-planes; any
extension to other domains is along specified analytic continuation.
Normalized primitives include their lower endpoint or their value at
the origin, so their additive constants are fixed.

All unrestricted statements are proved analytically or by exact finite
algebra in the article. The accompanying computations test conventions,
coefficients, and independent evaluation routes. Floating-point
agreement is reported as a diagnostic, never as a proof or as an
interval-certified error bound.

% END sections/01_scope.tex


% BEGIN sections/02_hurwitz.tex
\section{An entire normalization for all independent harmonic exponents}
\label{eht:sec}

The independent-inner-exponent question in the incoming exact-identities
report connects directly to an established interpolant of truncated
multiple zeta functions.  We identify that interpolant exactly, give a
self-contained residue proof of its entire continuation, and develop a
relative version with two endpoints.  The resulting endpoint, Gamma, and
Stieltjes identities give an analytic replacement for the one-parameter
Gamma normalization when the inner exponents move independently.

The result does not assert a reduction of every ordered coefficient to
ordinary zeta values.  In particular the nonsymmetric harmonic Stieltjes
coordinate from the incoming ordered-resonance article remains visible.
The analytic continuation and residue structure of multiple Hurwitz zeta
functions, and the quasi-shuffle algebra, are classical
\cite{eht:MehtaViswanadham,eht:Hoffman}.  The one-endpoint entire
interpolation and its stuffle law are due to Ihara, Nakamura, and Yamamoto,
with earlier overlapping interpolation work credited in their paper
\cite{eht:INY}.  The relative ordered quotient, residue calculation, and
explicit consequences below organize this known construction for the
repository's present questions.  No literature-wide priority claim is made.

\subsection{Definitions and the finite formula}

For $\Re a>0$ and a word $w=(s_1,\ldots,s_r)$, let
\begin{equation}\label{eht:Z}
 Z_a(w)=\sum_{n_1>\cdots>n_r\ge0}
              \prod_{j=1}^r(n_j+a)^{-s_j},\qquad Z_a(\varnothing)=1.
\end{equation}
Use principal powers.  The series is initially normally convergent when
$\Re(s_1+\cdots+s_j)>j$ for every $j$.  Its meromorphic continuation is
understood below.  The deliberately smaller domain $\Re s_j>1$ for every
$j$ suffices for all initial product manipulations.

For $\Re a,\Re b>0$, define recursively
\begin{align}
 \mathcal H_{a,b}(\varnothing)&=1,\notag\\
 \mathcal H_{a,b}(s_1,\ldots,s_r)
 &=Z_b(s_1,\ldots,s_r)
   -\sum_{j=1}^rZ_a(s_1,\ldots,s_j)
                    \mathcal H_{a,b}(s_{j+1},\ldots,s_r).
 \label{eht:rec}
\end{align}
The sum includes the final term $Z_a(s_1,\ldots,s_r)$.  Equivalently,
\begin{equation}\label{eht:finite}
 \mathcal H_{a,b}(w)=
 \sum_{q=0}^r(-1)^q
 \sum_{0=i_0<i_1<\cdots<i_q\le r}
 \left\{\prod_{\nu=1}^q
   Z_a(s_{i_{\nu-1}+1},\ldots,s_{i_\nu})\right\}
 Z_b(s_{i_q+1},\ldots,s_r).
\end{equation}
For $q=0$ the summand means $Z_b(w)$.  There are exactly $2^r$ terms,
counting repetitions before simplification.  Thus the definition requires
no limiting choice, ray, or finite-part prescription.

The first two cases are
\begin{align}
 h_{a,b}(s):=\mathcal H_{a,b}(s)
   &=\zeta(s,b)-\zeta(s,a),\label{eht:h}\\
 \mathcal H_{a,b}(s,t)
   &=Z_b(s,t)-Z_a(s,t)-\zeta(s,a)
                         \{\zeta(t,b)-\zeta(t,a)\}.
 \label{eht:depthtwo}
\end{align}
The second line already shows why subtracting just $Z_a-Z_b$ fails:
its residue at $s=1$ is generally nonzero.

There is also a shorter form in terms of the weakly ordered Hurwitz
function $Z_a^\star$, whose inequalities in \eqref{eht:Z} are all weak:
\begin{equation}\label{eht:star}
 \mathcal H_{a,b}(w)=\sum_{j=0}^r(-1)^j
       Z_a^\star(s_j,\ldots,s_1)Z_b(s_{j+1},\ldots,s_r).
\end{equation}
The empty prefix has value one.  The inverse of the strict ordered series
has coefficient $(-1)^jZ_a^\star(s_j,\ldots,s_1)$: expand each inverse
one-index factor, then reverse the weak order.  This proves
\eqref{eht:star} in an absolute-convergence domain and hence meromorphically.
Proposition 2.1 of \cite{eht:INY} therefore gives the exact identification
\begin{equation}\label{eht:INY}
                 \mathcal H_{a,1}(w)=\Psi(w;a-1).
\end{equation}
Their Corollary 4.2 already proves that this is entire in all indices;
the proof below provides a different route through explicit prefix
residues and treats the relative endpoints together.

\begin{theorem}[Entire relative harmonic interpolation]\label{eht:main}
For every depth $r$, the expression \eqref{eht:finite} extends to an
entire function of $(s_1,\ldots,s_r)\in\mathbb C^r$.  It is jointly
holomorphic with $a,b$ on the right half-plane.  Every fixed mixed spectral
derivative is consequently an ordinary derivative of one entire function,
including where several polar hyperplanes of its separate summands meet.

If $N\ge0$ is an integer and $a=b+N$, then
\begin{equation}\label{eht:interval}
 \boxed{\mathcal H_{b+N,b}(s_1,\ldots,s_r)=
  \sum_{N>n_1>\cdots>n_r\ge0}
              \prod_{j=1}^r(n_j+b)^{-s_j}.}
\end{equation}
In particular $\mathcal H_{M,1}$ is the strict finite multiple harmonic
sum with largest index at most $M-1$.
\end{theorem}

\subsection{A shift-independent prefix-residue lemma}

Write $S_j=s_1+\cdots+s_j$.  Introduce the rational-polynomial functions
\begin{equation}\label{eht:beta}
 \beta_0(z)=\frac1{z-1},\qquad
 \beta_1(z)=-\frac12,\qquad
 \beta_{2q}(z)=\frac{B_{2q}}{(2q)!}(z)_{2q-1}\quad(q\ge1),
 \qquad \beta_{2q+1}(z)=0\quad(q\ge1).
\end{equation}
Here the $B_{2q}$ are Bernoulli numbers, with $B_1=-1/2$.

\begin{lemma}[Prefix residue factorization]\label{eht:residue}
The possible polar divisors of $Z_a(s_1,\ldots,s_r)$ are among
\[
 s_1=1,\qquad S_k=k-m\quad(2\le k\le r,\ m\in\mathbb Z_{\ge0}).
\]
At a generic point of each divisor the pole is at most simple.  In the
normal coordinate $S_k-k+m$, its residue is
\begin{equation}\label{eht:resfactor}
 C_{k,m}(s_1,\ldots,s_k)Z_a(s_{k+1},\ldots,s_r),
\end{equation}
where $C_{1,0}=1$ and, for $k\ge2$,
\begin{equation}\label{eht:C}
 C_{k,m}=
 \sum_{\substack{e_1+\cdots+e_{k-1}=m\\e_i\ge0}}
 \prod_{i=1}^{k-1}
 \beta_{e_i}\left(S_i-(i-1)+\sum_{j<i}e_j\right),
 \qquad\text{restricted to }S_k=k-m.
\end{equation}
The coefficient is independent of $a$.  A vanishing $C_{k,m}$ simply means
that the indicated possible divisor is absent.
\end{lemma}

\begin{proof}
Sum first over $n_1$, put $x=n_2+a$, and use a finite
Euler--Maclaurin expansion
\begin{equation}\label{eht:EM}
 \zeta(z,x+1)=\frac{x^{1-z}}{z-1}-\frac12x^{-z}
    +\sum_{q=1}^Q\frac{B_{2q}}{(2q)!}(z)_{2q-1}x^{-z-2q+1}
    +R_Q(z,x).
\end{equation}
On any prescribed compact set of $z$, sufficiently large $Q$ makes the
remainder holomorphic there, with a uniform estimate
$R_Q(z,n+a)=O((n+1)^{1-\Re z-2Q})$ on compact right-half-plane sets of
$a$.  The same is true after fixed parameter derivatives, with logarithmic
factors.  This follows directly from the integral remainder with a bounded
periodic Bernoulli function.  In the remaining $r-1$ sums the number and
size of the inner terms have at most polynomial growth in $n_2$ on a
compact set of all exponents.  Taking $Q$ larger therefore makes the
remainder normally summable on that compact.

Each explicit term in \eqref{eht:EM} replaces the first two exponents by
$s_1+s_2-1+e$, with coefficient $\beta_e(s_1)$.  Repeat this step in the
lower-depth functions.  This proves meromorphic continuation by induction
and shows that every possible polar equation is a prefix equation of the
displayed form.  The pole at $s_1=1$ has residue
$Z_a(s_2,\ldots,s_r)$.

To reach a pole involving exactly the first $k$ indices, merge those
indices in $k-1$ steps.  Before step $i$ the current leading exponent is
$S_i-(i-1)+e_1+\cdots+e_{i-1}$.  After all merges, the final Hurwitz pole
has normal coordinate
$S_k-k+e_1+\cdots+e_{k-1}$.  The paths producing $S_k=k-m$ are exactly
the finite set in \eqref{eht:C}.  The remaining indices form the factor
$Z_a(s_{k+1},\ldots,s_r)$.  All coefficients in this procedure are
independent of the common shift.

The only denominators in the merging coefficients arise from proper
prefix hyperplanes.  Away from intersections, the final pole is simple.
Since $Q$ can be chosen arbitrarily large on a prescribed compact, the
holomorphic remainders cannot add further local poles.  This proves the
lemma, including its joint parameter assertion.
\end{proof}

For example,
\[
 C_{2,0}=\frac1{s_1-1},\qquad C_{2,1}=-\frac12,
 \qquad C_{2,2}=\frac{s_1}{12},\qquad C_{2,3}=0,
 \qquad C_{3,0}=\frac1{(s_1-1)(s_1+s_2-2)}.
\]
These are residues in a sum-of-indices coordinate; replacing that
coordinate by a scaled parameter rescales the residue.

\begin{proof}[Proof of Theorem~\ref{eht:main}]
Induct on the depth.  At depth one the common residue of the two Hurwitz
functions cancels.  Suppose all lower-depth $\mathcal H_{a,b}$ are entire.
Equation~\eqref{eht:rec} shows that the only remaining possible poles are
global prefix hyperplanes of the full word.  Fix one involving the first
$k$ indices and work away from all other hyperplanes.  The terms with
$j<k$ in \eqref{eht:rec} have no pole there.  By
Lemma~\ref{eht:residue}, the residue of the rest is
\[
 C_{k,m}\left\{
 Z_b(s_{k+1},\ldots,s_r)
 -\sum_{j=k}^rZ_a(s_{k+1},\ldots,s_j)
                  \mathcal H_{a,b}(s_{j+1},\ldots,s_r)
 \right\}=0.
\]
The bracket is exactly the lower-depth convolution identity
\eqref{eht:rec}, including the empty-prefix term at $j=k$.
Thus every generic polar hypersurface is removable.  The locally finite
union of their intersections has complex codimension at least two.
Hartogs extension removes these remaining intersections; equivalently,
one can cancel each linear factor in a local meromorphic denominator.
This proves joint entire continuation in the exponents.  The same proof
with $a,b$ in compact right-half-plane sets gives joint holomorphy.

For the interpolation assertion, take $\Re s_j>1$.  Split the strictly
ordered indices for $Z_b(w)$ at $N$: a uniquely determined initial block
lies in $n\ge N$, and the remaining block lies in $n<N$.  Hence
\[
 Z_b(w)=\sum_{j=0}^rZ_{b+N}(s_1,\ldots,s_j)
   \sum_{N>n_{j+1}>\cdots>n_r\ge0}
                    \prod_{\nu=j+1}^r(n_\nu+b)^{-s_\nu}.
\]
The triangular recursion identifies the finite sums with
$\mathcal H_{b+N,b}$.  Both sides are entire in all exponents, so the
identity extends to $\mathbb C^r$.
\end{proof}

\subsection{Composition, the stuffle product, and endpoint singularities}

For bookkeeping, attach a noncommuting letter $[s]$ to each exponent and
form the coefficientwise series
\[
 F_a=1+\sum_{r\ge1}\sum_{s_1,\ldots,s_r}
             Z_a(s_1,\ldots,s_r)[s_1]\cdots[s_r].
\]
This notation requires only a finite generating alphabet and its additive
closure at any application; it is not a sum over an uncountable set.
Concatenation of words gives
\begin{equation}\label{eht:quotient}
 \mathcal H_{a,b}=F_a^{-1}F_b.
\end{equation}

\begin{corollary}[Exact composition and shifts]\label{eht:transport}
With products interpreted by deconcatenation,
\begin{equation}\label{eht:cocycle}
 \boxed{\mathcal H_{a,b}(w)=
 \sum_{j=0}^r\mathcal H_{a,c}(s_1,\ldots,s_j)
                  \mathcal H_{c,b}(s_{j+1},\ldots,s_r).}
\end{equation}
Moreover
\begin{align}
 \mathcal H_{a+1,b}(w)
   &=\mathcal H_{a,b}(w)+a^{-s_1}
                              \mathcal H_{a,b}(s_2,\ldots,s_r),
       \label{eht:upper-shift}\\
 \mathcal H_{a,b}(w)
   &=\mathcal H_{a,b+1}(w)+b^{-s_r}
                              \mathcal H_{a,b+1}(s_1,\ldots,s_{r-1}).
       \label{eht:lower-shift}
\end{align}
All of these identities hold for arbitrary complex exponents.
\end{corollary}

\begin{proof}
The composition formula is
$F_a^{-1}F_cF_c^{-1}F_b=F_a^{-1}F_b$.
Splitting off the smallest index gives
$F_a=F_{a+1}E(a)$, where $E(a)=1+\sum_s a^{-s}[s]$.
Thus $\mathcal H_{a+1,b}=E(a)\mathcal H_{a,b}$ and
$\mathcal H_{a,b}=\mathcal H_{a,b+1}E(b)$.  Taking coefficients proves
the two shift formulas.
\end{proof}

The usual strict stuffle identity also holds without exclusions:
\begin{equation}\label{eht:stuffle-two}
 h_{a,b}(s)h_{a,b}(t)=
 \mathcal H_{a,b}(s,t)+\mathcal H_{a,b}(t,s)+h_{a,b}(s+t).
\end{equation}
More generally $\mathcal H_{a,b}$ is a character of the quasi-shuffle
algebra whose merged-letter operation is $[s]\circ[t]=[s+t]$.  Indeed $Z_a$ and
$Z_b$ are characters in the initial domain by splitting equality and
inequality cases in a product of sums.  This remains true meromorphically.
Characters into a commutative algebra are closed under convolution and
convolution inversion in the connected quasi-shuffle Hopf algebra
\cite{eht:Hoffman}; \eqref{eht:quotient} therefore gives the assertion.
Its coefficients are entire by Theorem~\ref{eht:main}.

Inverting the first shift formula gives the useful finite identity
\begin{equation}\label{eht:inverse-shift}
 \boxed{\mathcal H_{a,b}(s_1,\ldots,s_r)=
  \sum_{j=0}^r(-1)^j a^{-S_j}
                  \mathcal H_{a+1,b}(s_{j+1},\ldots,s_r),}
 \qquad S_0=0.
\end{equation}
It continues the upper endpoint to the principal slit plane
$\mathbb C\setminus(-\infty,0]$, by shifting it finitely many times into
the right half-plane.  Near $a=0$ the right side is a finite sum of
$a^{-S_j}$ times functions holomorphic in $a$.  Taylor expansion of those
functions gives a \emph{convergent} local expansion.  Thus the exact
branch exponents are prefix sums, and every fixed spectral derivative
adds only finitely many powers of $\log a$.

Continuation near $a=-N$, on a specified local branch, follows from
\[
 \mathcal H_{a,b}=E(a)^{-1}E(a+1)^{-1}\cdots E(a+N)^{-1}
                         \mathcal H_{a+N+1,b}.
\]
At fixed depth this is finite coefficient algebra.  Only the factor
$E(a+N)^{-1}$ is locally singular, so the branch exponents are sums over
consecutive subwords.  No essential singularity is introduced.  The
lower endpoint has the simpler local form \eqref{eht:lower-shift}: near
$b=0$ it has one possible singular factor $b^{-s_r}$.

\subsection{Gamma closure of the symmetric integer-index sector}

The symmetric-sum theorem for a stuffle character gives
\begin{equation}\label{eht:symmetric}
 \sum_{\sigma\in S_r}\mathcal H_{a,b}(s_{\sigma(1)},\ldots,s_{\sigma(r)})
 =\sum_{\pi}(-1)^{r-|\pi|}
       \prod_{B\in\pi}(|B|-1)!\,
       \prod_{B\in\pi}h_{a,b}\left(\sum_{j\in B}s_j\right),
\end{equation}
where $\pi$ runs over set partitions of $\{1,\ldots,r\}$.
This is the classical symmetric-sum mechanism applied to the entire
character.  Repeated labels are counted with their permutation
multiplicities; no independence of zeta values is assumed.

In particular, as a convergent germ at $u=0$,
\begin{equation}\label{eht:ones-Gamma}
 \boxed{\sum_{r\ge0}\mathcal H_{a,b}(\{1\}^r)u^r
       =\frac{\Gamma(a+u)\Gamma(b)}{\Gamma(a)\Gamma(b+u)}.}
\end{equation}
Its first coefficients, with $h_j=h_{a,b}(j)$, are
\[
 \mathcal H_{a,b}(1,1)=\frac{h_1^2-h_2}{2},\qquad
 \mathcal H_{a,b}(1,1,1)=\frac{h_1^3-3h_1h_2+2h_3}{6},
 \qquad h_1=\psi(a)-\psi(b).
\]

There is a simultaneous finite Gamma formula for every finite alphabet
of positive integer exponents.  Let $p_1,\ldots,p_d\ge1$, let
$M=\max_jp_j$, and let $\rho_1,\ldots,\rho_M$ be the roots, with
multiplicity, of
\[
 x^M+\sum_{j=1}^du_jx^{M-p_j}.
\]
Then
\begin{equation}\label{eht:alphabet-Gamma}
 \boxed{\sum_{r\ge0}\sum_{j_1,\ldots,j_r=1}^d
  \mathcal H_{a,b}(p_{j_1},\ldots,p_{j_r})u_{j_1}\cdots u_{j_r}
  =\prod_{\ell=1}^M
       \frac{\Gamma(a-\rho_\ell)\Gamma(b)}
            {\Gamma(a)\Gamma(b-\rho_\ell)}.}
\end{equation}
The right side is symmetric in the roots and therefore a holomorphic
germ in the $u_j$, even at colliding roots.  It evaluates sums over
permutations of the prescribed indices; it does not identify individual
ordered words with each other.

To prove \eqref{eht:alphabet-Gamma}, abelianize the quasi-shuffle series.
Its logarithm is the formal difference of
$\sum_n\log(1+\sum_ju_j(n+b)^{-p_j})$ and the corresponding $a$ series.
This statement first follows when all independent indices have real part
greater than one; the symmetric-sum identity and the entire differences
$h_{a,b}$ continue it to the integer alphabet.  Each coefficient of the
difference converges, including the linear exponent-one coefficient.
Factoring the polynomial gives the logarithm
\[
 -\sum_{m\ge1}\frac{h_{a,b}(m)}m
                     \sum_{\ell=1}^M\rho_\ell^m.
\]
The Taylor expansion of
$\log\Gamma(a-\rho)-\log\Gamma(a)
 -\log\Gamma(b-\rho)+\log\Gamma(b)$ is exactly this expression.
This proves the identity and also supplies a branch-independent
coefficient algorithm.

\subsection{Independent logarithmic jets and the ordering coordinate}

Put $\gamma_m(x)$ for the generalized Stieltjes constants in
\[
 \zeta(1+u,x)=\frac1u+
       \sum_{m\ge0}\frac{(-1)^m\gamma_m(x)}{m!}u^m.
\]
Then
\begin{equation}\label{eht:h-jets}
 h_{a,b}^{(m)}(1)=(-1)^m\{\gamma_m(b)-\gamma_m(a)\}.
\end{equation}
Differentiating \eqref{eht:stuffle-two} on the diagonal gives the exact
first-jet identity
\begin{align}\label{eht:sum-jets}
 (\partial_1+\partial_2)\mathcal H_{a,b}(1,1)
  ={}&\{\psi(a)-\psi(b)\}\{\gamma_1(a)-\gamma_1(b)\}\notag\\
    &-\zeta'(2,b)+\zeta'(2,a).
\end{align}

The incoming ordered-resonance paper uses the coefficient $\eta(a)$ in
\[
 Z_a(1+u,1)=u^{-2}+\gamma_0(a)u^{-1}
       +\frac{\gamma_0(a)^2-\zeta(2,a)}2+\eta(a)u+O(u^2).
\]
Substitution in \eqref{eht:depthtwo} now gives
\begin{align}
 \partial_1\mathcal H_{a,b}(1,1)
   &=\eta(b)-\eta(a)+\gamma_1(a)\{\psi(a)-\psi(b)\},
       \label{eht:first-eta}\\
 \partial_2\mathcal H_{a,b}(1,1)
   &=\eta(a)-\eta(b)-\gamma_1(b)\{\psi(a)-\psi(b)\}
                     -\zeta'(2,b)+\zeta'(2,a).
       \label{eht:second-eta}
\end{align}
These formulas explain the scope of the cancellation theorem precisely.
It removes the moving poles and makes the separate logarithmic
decorations well defined; it does not remove the nonsymmetric arithmetic
coordinate $\eta$.

\subsection{Finite reductions with a nonpositive inner exponent}

One independent-index subfamily does have complete depth-one closure.
For $m\ge0$ and arbitrary complex $s$,
\begin{equation}\label{eht:negative-inner}
 \boxed{\mathcal H_{a,b}(s,-m)=\frac1{m+1}
 \left\{\sum_{j=0}^{m+1}\binom{m+1}{j}B_{m+1-j}\,
                       h_{a,b}(s-j)
          -B_{m+1}(b)h_{a,b}(s)\right\}.}
\end{equation}
All apparent singular values in the separate Hurwitz functions are
removable within the $h$ differences.  In particular,
\begin{align}
 \mathcal H_{a,b}(s,0)&=h_{a,b}(s-1)-b\,h_{a,b}(s),
        \label{eht:zero-inner}\\
 \mathcal H_{a,b}(0,s)&=(a-1)h_{a,b}(s)-h_{a,b}(s-1),\notag\\
 \mathcal H_{a,b}(s,-1)&=\frac12\{h_{a,b}(s-2)-h_{a,b}(s-1)
                                -b(b-1)h_{a,b}(s)\}.\notag
\end{align}

For a direct proof, first take $\Re s>m+2$ and use the ordinary finite
power-sum identity
\[
 \sum_{k=0}^{n-1}(k+a)^m
       =\frac{B_{m+1}(n+a)-B_{m+1}(a)}{m+1}
\]
in $Z_a(s,-m)$.  Expand $B_{m+1}(n+a)$ in powers of $n+a$.
Insert the resulting finite Hurwitz formula, and its $b$ version, into
\eqref{eht:depthtwo}.  Since
$h_{a,b}(-m)=\{B_{m+1}(a)-B_{m+1}(b)\}/(m+1)$, the terms involving
$B_{m+1}(a)$ cancel.  This yields \eqref{eht:negative-inner}.  Entire
continuation gives every $s$.  The second formula in the displayed list
also follows from the stuffle identity.

As explicit Gamma--Stieltjes specializations at $b=1$,
\begin{align}
 \mathcal H_{a,1}(1,0)&=a-1-\psi(a)-\gamma,\notag\\
 \partial_1\mathcal H_{a,1}(1,0)
   &=-\log\Gamma(a)-\gamma_1(a)+\gamma_1,\label{eht:Gamma-Stieltjes}\\
 \partial_2\mathcal H_{a,1}(0,1)
   &=(a-1)\{\gamma_1(a)-\gamma_1\}+\log\Gamma(a).\notag
\end{align}
The derivatives are taken before the indicated index specialization.
The identity $h_{a,1}'(0)=-\log\Gamma(a)$ fixes the normalization.

These reductions also have normalized antiderivatives.  For a path in
the right half-plane from $b$ to $a$, define
\begin{equation}\label{eht:primitive}
 P_{a,b}(s)=\int_b^a h_{x,b}(s)\,dx
    =(a-b)\zeta(s,b)
     +\frac{\zeta(s-1,a)-\zeta(s-1,b)}{s-1}.
\end{equation}
The combined right side is entire in $s$.  At $s=1$ it equals
\[
 P_{a,b}(1)=\log\Gamma(a)-\log\Gamma(b)-(a-b)\psi(b).
\]
Replacing each $h_{a,b}(s-j)$ in \eqref{eht:negative-inner} by
$P_{a,b}(s-j)$ gives the uniquely normalized integral
$\int_b^a\mathcal H_{x,b}(s,-m)\,dx$.  Every spectral derivative
commutes with this identity.  Thus this independent-index family has a
finite calculus in Hurwitz-zeta derivatives, Gamma functions, and
generalized Stieltjes constants, with its additive constant fixed.

\subsection{Polylogarithmic dictionary}

For $|z|<1$, the interpolation theorem gives
\begin{equation}\label{eht:polylog}
 \sum_{N\ge0}\mathcal H_{N+1,1}(s_1,\ldots,s_r)z^N
  =\frac{\operatorname{Li}_{s_1,\ldots,s_r}(z)}{1-z},
\end{equation}
where the strict multiple polylogarithm has numerator $z^{n_1}$.
This is an ordinary normally convergent identity, locally in all
independent exponents.  Each spectral derivative inserts the expected
logarithmic factors in the harmonic sums, and the normalized Euler
primitive of the numerator raises its first index by one.  These
classical correspondences now attach to an entire interpolation in the
finite harmonic endpoint.

This correspondence also explains the compatibility with the positive-word
shuffle calculus in Section~\ref{mixed:section}: finite harmonic endpoint
coefficients and nested polylogarithmic words encode the same strict
ordering. The independent-order kernel and nonsymmetric-jet problems
are collected in Section~\ref{research:section}.

% END sections/02_hurwitz.tex


% BEGIN sections/03_tornheim.tex
\section{All cubic Tornheim directions from two scalar coordinates}
\label{tr3:section}

The incoming \emph{Coincident Stieltjes and Directional Zeta} report~\cite{tr3:incoming} proves
all quadratic derivatives of $T(at,bt,ct)$ and asks for the entire cubic
layer.  The following calculation answers that question with two explicit
scalar coordinates.  One is the already isolated diagonal derivative
$\Omega$; the other is an absolutely convergent logarithmic Gamma series.
The result also identifies an exact conic of directions and a cyclic
symmetrization in which the second coordinate cancels.  These are finite
identities; no arithmetic independence of the two coordinates is assumed.

Throughout this section,
\[
 T(A,B,C)=\sum_{m,n\geq1}m^{-A}n^{-B}(m+n)^{-C},
 \qquad L=\log(2\pi),
\]
initially in its domain of absolute convergence and subsequently by its
meromorphic continuation.  Set
\[
 z_2=\zeta''(0),\qquad z_3=\zeta'''(0),\qquad
 y_3=\zeta'''(-1),\qquad
 \Omega=\left.\frac{d^3}{dt^3}T(t,t,t)\right|_{t=0}.
\]
The order of restriction is part of the notation.  In particular, $T$ is
not jointly holomorphic at the origin, whereas the one-variable diagonal
is holomorphic there~\cite{tr3:BD}.  The Stieltjes convention is
\[
 \zeta(1+u)=\frac1u+\sum_{r\geq0}\frac{(-1)^r\gamma_r}{r!}u^r.
\]

\subsection{A local remainder with no Gamma factors in its polar part}

\begin{lemma}[Elementary polar model]\label{tr3:polar}
There is a unique function $R(A,B,C)$ holomorphic in a neighborhood of
$(0,0,0)$, symmetric in $A,B$, for which
\begin{equation}\label{tr3:model}
 \boxed{\displaystyle
 T(A,B,C)=\zeta(A)\zeta(B)
  -\frac{C\zeta(B-1)}{A+C}
  -\frac{C\zeta(A-1)}{B+C}+C R(A,B,C).}
\end{equation}
The identity is meromorphic; each quotient retains its indicated
numerator until a restriction has been specified.
\end{lemma}

\begin{proof}
For $A,B$ near zero, Jonqui\`ere's expansion~\cite[\S25.12(ii)]{tr3:DLMF} gives
\[
 \operatorname{Li}_A(e^{-x})
 =\Gamma(1-A)x^{A-1}
   +\sum_{j\geq0}\zeta(A-j)\frac{(-x)^j}{j!},\qquad 0<x<2\pi.
\]
In the Mellin integral
\[
 \Gamma(C)T(A,B,C)=\int_0^\infty
 x^{C-1}\operatorname{Li}_A(e^{-x})\operatorname{Li}_B(e^{-x})\,dx,
\]
subtract on $(0,1)$ the six terms
\begin{align*}
 S(A,B;x)={}&\Gamma(1-A)\Gamma(1-B)x^{A+B-2}
 +\Gamma(1-A)\zeta(B)x^{A-1}
 +\Gamma(1-B)\zeta(A)x^{B-1}\\
 &-\Gamma(1-A)\zeta(B-1)x^A
 -\Gamma(1-B)\zeta(A-1)x^B+\zeta(A)\zeta(B).
\end{align*}
The difference is normally integrable against $x^{C-1}$ on a sufficiently
small common parameter polydisc.  Each fixed parameter derivative only
adds logarithmic powers.  Integrating the six displayed monomials shows
that some holomorphic symmetric $H$ satisfies
\begin{align}\label{tr3:Hmodel}
 T(A,B,C)=\frac{1}{\Gamma(1+C)}\biggl\{
 \zeta(A)\zeta(B)
 -\frac{C\Gamma(1-A)\zeta(B-1)}{A+C}
 -\frac{C\Gamma(1-B)\zeta(A-1)}{B+C}+C H(A,B,C)\biggr\}.
\end{align}
The terms with denominators $A+B+C-2$, $A+C-1$, and $B+C-1$ belong
to $H$, since these denominators are nonzero at the origin.

Write $G(C)=1/\Gamma(1+C)$.  Subtract the proposed polar part from
\eqref{tr3:Hmodel} and divide by $C$.  The resulting expression is
\begin{align*}
 R(A,B,C)={}&G(C)H(A,B,C)
  +\frac{G(C)-1}{C}\zeta(A)\zeta(B)\\
 &-\frac{\Gamma(1-A)G(C)-1}{A+C}\zeta(B-1)
  -\frac{\Gamma(1-B)G(C)-1}{B+C}\zeta(A-1).
\end{align*}
Every quotient is holomorphic.  For example,
$\Gamma(1-A)G(-A)=1$, so its numerator is divisible by $A+C$ in
the local ring of holomorphic functions.  This proves existence and
symmetry.  Uniqueness follows from \eqref{tr3:model} for $C\ne0$ and
then by analytic continuation.
\end{proof}

The meromorphic continuation and Mellin method are classical~\cite{tr3:BD,tr3:SS}; the local
form \eqref{tr3:Hmodel} also occurs in the incoming report.  The change
from $H$ to $R$ is useful because all Gamma and Euler constant terms
are absorbed into a holomorphic function before the coefficient
calculation.

For reference, the complete linear jet is
\begin{align}\label{tr3:linearjet}
 R(0,0,0)&=\frac L2+\zeta'(-1),\notag\\
 R_A(0)=R_B(0)&=\frac{L^2}{8}-\frac{z_2}{2},\qquad
 R_C(0)=\frac{\zeta''(-1)-z_2}{2}.
\end{align}
These formulas need no additional spectral input.  The exact $C$ axis
in the next proof determines $R(0)$ and $R_C(0)$.  The coefficient of
$AC$ in the Euler identity \eqref{tr3:stuffleR} gives
$2R_A=\zeta'(0)^2-\zeta''(0)$.  As a consequence,
\[
 \left.\frac{d^2}{dt^2}T(t,t,t)\right|_{t=0}
  =\frac{L^2}{2}-z_2-\zeta''(-1)+4R_A+2R_C
  =L^2-4z_2.
\]
This recovers the incoming quadratic diagonal formula by an independent
coefficient calculation, while the new result begins at cubic order.

\subsection{The quadratic remainder and the cubic ray theorem}

Let all partial derivatives in the following notation be evaluated
at $(0,0,0)$:
\[
 p=R_{AA}=R_{BB},\qquad q=R_{AB},\qquad
 r=R_{AC}=R_{BC},\qquad s=R_{CC}.
\]
These four coordinates describe the homogeneous quadratic part of
$R$.  Three exact restrictions determine three independent combinations.

\begin{lemma}[Three exact constraints]\label{tr3:constraints}
The four second partial derivatives satisfy
\begin{equation}\label{tr3:quadraticjet}
 \boxed{\displaystyle
 s=\frac{y_3-z_3}{3},\qquad
 p+2r=-\frac{Lz_2}{2}-z_3,\qquad
 q=\frac{\Omega+6Lz_2+8z_3}{6}.}
\end{equation}
\end{lemma}

\begin{proof}
Counting the $n-1$ positive pairs with sum $n$ gives
\[
 T(0,0,C)=\zeta(C-1)-\zeta(C).
\]
Since $\zeta(0)^2-2\zeta(-1)=5/12$, \eqref{tr3:model} yields
\[
 C R(0,0,C)=\zeta(C-1)-\zeta(C)-\frac5{12}.
\]
Its cubic coefficient gives $s=(y_3-z_3)/3$.

On the plane $B=0$, set
\[
 E(A,C)=T(A,0,C)=\sum_{n>m\geq1}n^{-C}m^{-A}.
\]
The ordinary Euler stuffle identity, first in an absolutely convergent
region and then meromorphically, is
\[
 E(A,C)+E(C,A)=\zeta(A)\zeta(C)-\zeta(A+C).
\]
Inserting \eqref{tr3:model} gives a holomorphic identity
\begin{align}\label{tr3:stuffleR}
 C R(A,0,C)+A R(C,0,A)
 ={}&\zeta(A)\zeta(C)-\zeta(A+C)
 +\frac{\zeta(A)+\zeta(C)}2\notag\\
 &+\zeta(-1)+\zeta(A-1)+\zeta(C-1).
\end{align}
The coefficient of $A^2C$ on its left is $p/2+r$.  On its right it is
$\zeta'(0)z_2/2-z_3/2=-Lz_2/4-z_3/2$.  This proves the second
constraint.

Finally, on the diagonal \eqref{tr3:model} says
\[
 T(t,t,t)=\zeta(t)^2-\zeta(t-1)+tR(t,t,t).
\]
Taking third derivatives gives
\[
 \Omega=-3Lz_2-z_3-y_3+6p+6q+12r+3s.
\]
The first two constraints simplify this to
$\Omega=-6Lz_2-8z_3+6q$, proving the last assertion.
\end{proof}

\begin{theorem}[Every asymmetric cubic direction]\label{tr3:main}
Let $a,b,c\in\mathbb C$ satisfy $(a+c)(b+c)\ne0$.  Form the
one-variable meromorphic restriction $W_{a,b,c}(t)=T(at,bt,ct)$ before
continuing to $t=0$.  This restriction is holomorphic at zero.  Define
$\kappa=R_{AA}(0,0,0)$.  Then
\begin{align}\label{tr3:allrays}
 \boxed{\displaystyle W_{a,b,c}'''(0)}
 ={}&abc\,\Omega
 +3c\bigl(a^2+b^2-c(a+b)\bigr)\kappa\notag\\
 &-\frac32\bigl(ab(a+b-4c)+c^2(a+b)\bigr)Lz_2\notag\\
 &+\left(-\frac{a^3+b^3}{2}+8abc-3c^2(a+b)-c^3\right)z_3\notag\\
 &+\left(c^3-\frac{cb^3}{a+c}-\frac{ca^3}{b+c}\right)y_3.
\end{align}
The scalar $\kappa$ is explicitly given by the absolutely convergent
series in \eqref{tr3:kappa} below.
\end{theorem}

\begin{proof}
On an admissible ray the two polar quotients in \eqref{tr3:model}
become constants times holomorphic one-variable zeta functions, so
the restriction is holomorphic.  Its cubic derivative is
\begin{align*}
 W_{a,b,c}'''(0)={}&-\frac{a^3+b^3}{2}z_3
 -\frac32ab(a+b)Lz_2
 -c\left(\frac{b^3}{a+c}+\frac{a^3}{b+c}\right)y_3\\
 &+3cp(a^2+b^2)+6cqab+6c^2r(a+b)+3c^3s.
\end{align*}
Substitute $p=\kappa$ and \eqref{tr3:quadraticjet}.  Collecting the
coefficients of $\Omega,\kappa,Lz_2,z_3,y_3$ gives
\eqref{tr3:allrays}.
\end{proof}

Thus no unresolved partial derivative of the three-variable Tornheim
function remains in the answer.  The new information relative to the
quadratic theorem is one explicitly specified Gamma-series coordinate
in addition to the diagonal coordinate already present in the cubic
digamma problem.  This is a two-coordinate reduction, not a claim that
two algebraically independent constants are necessary.

\subsection{A convergent Gamma series and all boundary jets}

Put
\begin{equation}\label{tr3:rho}
 \rho(n)=\log\Gamma(n+1)-\left(n+\frac12\right)\log n
            +n-\frac L2-\frac{1}{12n},\qquad n\geq1.
\end{equation}
Stirling's expansion~\cite[\S5.11]{tr3:DLMF} gives $\rho(n)=-1/(360n^3)+O(n^{-5})$.
Consequently
\[
 M_j=\sum_{n\geq1}(\log n)^j\rho(n)
\]
converges absolutely for every fixed integer $j\geq0$.  No zeta
regularization is used in these sums.

\begin{theorem}[Boundary germ in Gamma and Stieltjes data]
\label{tr3:boundary}
The Dirichlet series
\[
 D(A)=\sum_{n\geq1}\frac{\rho(n)}{n^A}
\]
is holomorphic on $\Re A>-2$.  Near $A=0$ it satisfies
\begin{align}\label{tr3:boundarygerm}
 R(A,0,0)={}&D(A)-\zeta'(A-1)-\frac12\zeta'(A)
            -\zeta(A-1)
            +\frac{\zeta(1+A)-1/A}{12}.
\end{align}
In particular, for every $j\geq0$,
\begin{align}\label{tr3:boundaryjets}
 \partial_A^jR(0,0,0)
 ={}&(-1)^j M_j-\zeta^{(j+1)}(-1)
 -\frac12\zeta^{(j+1)}(0)-\zeta^{(j)}(-1)
 +\frac{(-1)^j\gamma_j}{12},
\end{align}
and the scalar in Theorem~\ref{tr3:main} is
\begin{equation}\label{tr3:kappa}
 \boxed{\displaystyle
 \kappa=\sum_{n\geq1}(\log n)^2\rho(n)
       -\zeta'''(-1)-\frac12\zeta'''(0)-\zeta''(-1)
       +\frac{\gamma_2}{12}.}
\end{equation}
\end{theorem}

\begin{proof}
The decay of $\rho$ proves normal convergence of $D$ and every fixed
order derivative on compact subsets of $\Re A>-2$.  For $\Re A>2$
and $C$ near zero, the representation
\[
 E(A,C)=\sum_{m\geq1}m^{-A}\zeta(C,m+1)
\]
can be differentiated normally in $C$.  Lerch's formula
$\zeta'(0,x)=\log\Gamma(x)-L/2$~\cite[(25.11.18)]{tr3:DLMF} gives
\[
 E_C(A,0)=\sum_{m\geq1}m^{-A}
            \left(\log\Gamma(m+1)-\frac L2\right).
\]
Expanding the summand using \eqref{tr3:rho}, in this convergent
half-plane, proves
\[
 E_C(A,0)=D(A)-\zeta'(A-1)-\frac12\zeta'(A)
          -\zeta(A-1)+\frac{\zeta(A+1)}{12}.
\]
Both sides have meromorphic continuations to a neighborhood of zero.
On the other hand, differentiation in $C$ of \eqref{tr3:model} with
$B=0$, for fixed $A\ne0$, gives
\[
 E_C(A,0)=\frac{1}{12A}+R(A,0,0).
\]
Combining these two identities proves \eqref{tr3:boundarygerm}; its
apparent $A=0$ singularity is removable.  Taking $j$ derivatives and
using the specified Stieltjes convention proves
\eqref{tr3:boundaryjets} and \eqref{tr3:kappa}.
\end{proof}

The first two members of this hierarchy are completely evaluable:
\begin{align}\label{tr3:lowGamma}
 \sum_{n\geq1}\rho(n)
   &=2\zeta'(-1)+\frac L4-\frac{1+\gamma}{12},\notag\\
 \sum_{n\geq1}(\log n)\rho(n)
   &=-\frac{L^2}{8}-\zeta''(-1)-\zeta'(-1)-\frac{\gamma_1}{12}.
\end{align}
They follow from \eqref{tr3:linearjet} and
\eqref{tr3:boundaryjets}.  We use these as derived consistency
identities and make no separate novelty claim for them.

There is also an exact integral involving an order derivative of the
ordinary polylogarithm.  In the half-plane $\Re A>-2$,
\begin{equation}\label{tr3:Dpolylog}
 D(A)=\int_0^\infty
 \left(\frac1{e^x-1}-\frac1x+\frac12-\frac{x}{12}\right)
           \frac{\operatorname{Li}_A(e^{-x})}{x}\,dx.
\end{equation}
To see this, the Laplace form of Binet's formula~\cite[\S5.9]{tr3:DLMF} gives
\[
 \rho(n)=\int_0^\infty e^{-nx}
 \left(\frac1{e^x-1}-\frac1x+\frac12-\frac{x}{12}\right)\frac{dx}{x}.
\]
This formula can also be verified by differentiating with respect to
$n$ using the standard digamma integral and fixing the constant at
$n=+\infty$.  Summation and integration commute absolutely on
$\Re A>-2$.  At the lower endpoint the bracket divided by $x$ is
$-x^2/720+O(x^4)$ and the polylogarithm contributes at worst the
integrable singularity required by that half-plane; the upper endpoint
is exponentially decaying.  Fixed order derivatives add logarithmic
factors and remain integrable.  Hence
\begin{equation}\label{tr3:Mpolylog}
 M_j=(-1)^j\int_0^\infty
 \left(\frac1{e^x-1}-\frac1x+\frac12-\frac{x}{12}\right)
 \left.\partial_A^j\operatorname{Li}_A(e^{-x})\right|_{A=0}\frac{dx}{x}.
\end{equation}
Equations \eqref{tr3:kappa} and \eqref{tr3:Mpolylog} supply equivalent
Gamma-series and polylogarithmic-integral coordinates for the same
constant.  They also connect the entire boundary derivative hierarchy
to the Stieltjes constants.  Since
$-\left.\partial_s H_n^{(s)}\right|_{s=0}=\log\Gamma(n+1)$,
these identities can equivalently be written with order derivatives of
generalized harmonic numbers.

\subsection{Exact cancellation families and examples}

\begin{corollary}[A conic of directions]\label{tr3:conic}
If an admissible direction obeys
\begin{equation}\label{tr3:coniceq}
 a^2+b^2=c(a+b),
\end{equation}
then its third derivative is a finite combination of
$\Omega$, $L\zeta''(0)$, $\zeta'''(0)$, and $\zeta'''(-1)$.
No Gamma-series coordinate occurs.
\end{corollary}

For example, the direction $(1,2,5/3)$ gives the exact relation
\begin{equation}\label{tr3:conicexample}
 W_{1,2,5/3}'''(0)=\frac{10}{3}\Omega
 -\frac32Lz_2-\frac{403}{54}z_3-\frac{245}{297}y_3.
\end{equation}
The diagonal and the Euler diagonal are familiar points on the same
conic.  In particular,
\[
 W_{1,0,1}'''(0)=-\frac32Lz_2-\frac92z_3,
\]
as also follows directly from
$T(t,0,t)=(\zeta(t)^2-\zeta(2t))/2$.

\begin{corollary}[Cyclic elimination]\label{tr3:cyclic}
Let all pairwise sums $a+b$, $a+c$, $b+c$ be nonzero, and put
$e_1=a+b+c$, $e_2=ab+bc+ca$, $e_3=abc$.  Then
\begin{align}\label{tr3:cyclicformula}
 W_{a,b,c}'''(0)+W_{b,c,a}'''(0)+W_{c,a,b}'''(0)
 ={}&3e_3\Omega+3(9e_3-e_1e_2)Lz_2\notag\\
 &+(-2e_1^3+3e_1e_2+27e_3)z_3.
\end{align}
Both $\kappa$ and $\zeta'''(-1)$ cancel.
\end{corollary}

\begin{proof}
Sum \eqref{tr3:allrays} cyclically.  The $\kappa$ coefficient is
$3\sum_{\rm cyc}c(a^2+b^2)-3\sum_{\rm cyc}c^2(a+b)=0$.
In the $y_3$ coefficient, terms sharing the denominator $a+c$ combine
to $-b^3$, cancelling the sum of cubes.  The remaining two
coefficients reduce to the elementary symmetric expressions displayed
in \eqref{tr3:cyclicformula}.
\end{proof}

The concrete specialization $(a,b,c)=(1,2,3)$ is
\[
 W_{1,2,3}'''(0)+W_{2,3,1}'''(0)+W_{3,1,2}'''(0)-18\Omega
       =-36Lz_2-72z_3.
\]
Individually, two asymmetric cases read
\begin{align*}
 W_{1,2,3}'''(0)&=6\Omega-36\kappa
                -\frac{27}{2}Lz_2-\frac{129}{2}z_3+\frac{102}{5}y_3,\\
 W_{2,3,1}'''(0)&=6\Omega+24\kappa
                -\frac{33}{2}Lz_2+\frac{29}{2}z_3-10y_3.
\end{align*}
The Euler direction $(1,0,2)$ isolates the additional coordinate:
\begin{equation}\label{tr3:kappaisolate}
 W_{1,0,2}'''(0)=-6\kappa-6Lz_2-\frac{41}{2}z_3+7y_3.
\end{equation}
Thus one may replace $\kappa$ by a single Euler double-zeta directional
derivative if that basis is preferable.  The Gamma series gives a
convergent evaluation of that derivative, rather than merely naming it.

For a finite linear combination of admissible rays, the map
\[
 (a,b,c)\longmapsto
 \bigl(abc,\ 3c(a^2+b^2-c(a+b))\bigr)
\]
lists the coefficients of the two residual coordinates.  Any linear
combination in the kernel of this explicit map reduces to the displayed
ordinary zeta derivatives.  This is a sufficient exact elimination
criterion.  Arithmetic relations between $\Omega$ and $\kappa$, if
found later, could enlarge the actual reducible family.

\subsection{Verification and status}

The accompanying \texttt{check\_cubic\_rays.py} checks the finite algebra
symbolically.  It computes the Gamma series using finite Stirling
subtraction and evaluates Tornheim functions independently through
integrated Jonqui\`ere series on $(0,1)$ and an incomplete-Gamma series
on $(1,\infty)$.  The reported working precision, finite-difference
step, truncations, and discrepancies are recorded in
\texttt{cubic\_rays.json}.  The largest discrepancy in the three fully asymmetric cases is less than
$2.27\times10^{-38}$ at 65-digit working precision.  These numerical
diagnostics accompany the proof; the floating-point results are not interval certificates.
For orientation,
\begin{align*}
 \kappa&=3.6336347857639257352355574698800935822\ldots,\\
 M_2&=-0.00064921643050147820901953512402746385\ldots.
\end{align*}

The exact contribution is the cubic directional classification,
the explicit boundary germ, and their cancellation consequences.
The classical ingredients are Mellin continuation, Euler stuffle,
Stirling subtraction, and Lerch's formula.  The 2025 preprint of
Sathyanarayana and Sharan~\cite{tr3:SS}, published in 2026, develops fixed-variable
Kronecker limits, Herglotz-type relations, and diagonal special values.
Its Section~5 proposes further study of simultaneously varying order
limits at points of indeterminacy.  The 2026 preprint of
Dobrowolski~\cite{tr3:Dobrowolski} concerns positive-integer symmetrized weights.  These are
relevant adjacent developments.  The simultaneous cubic-origin formula
above was not located in the inspected versions, but a comprehensive
priority claim would require a broader literature review.

Questions about further reduction of $\Omega$ and $\kappa$, the quartic
jet, and systematic cancellation families are stated in
Section~\ref{research:section}.

% END sections/03_tornheim.tex


% BEGIN sections/04_mixed.tex
\section{A finite reduction for arbitrary mixed integer inner orders}
\label{mixed:section}

The directional-zeta report reduces two integer inner orders and the
arbitrary-depth sector in which every inner order is nonpositive. Its
Question 11 asks for a finite nested-polylogarithm reduction at arbitrary
depth with both signs present, while the outer order stays complex
\cite{RepoDirectional}. We solve that formulation. The construction uses
root filtering, shuffle products, rational partial fractions, and finite
geometric power sums. Each stage is explicit and accommodates repeated
colors. The analytic continuation and the directions in which one may
differentiate are part of the statement.

The shuffle algebra and rational manipulation of hyperlogarithms are
classical and have powerful general implementations \cite{Panzer}.
Here the desired output is a finite expression with one free complex
order, positive remaining indices, a controlled depth, and the fixed
cumulative root alphabet. The construction and its proof are given in
full, so their correctness does not depend on an external symbolic
integration package.

\subsection{Statement, alphabet, and convergence}

For integers $a_j$, positive integers $p_j$, and nonzero colors
$|X_j|\le1$, put
\begin{equation}\label{mixed:Z}
 Z(C)=\sum_{m_1,\ldots,m_d\ge1}
       \frac{\prod_{j=1}^d X_j^{m_j}m_j^{-a_j}}
            {(p_1m_1+\cdots+p_dm_d)^C}.
\end{equation}
If all colors have modulus less than one, this series is entire in
$C$. For unit colors it initially converges absolutely when, for example,
\[
              \Re C>d+\sum_j\max(0,-a_j).
\]
This sufficient half-plane follows by grouping according to the
denominator and using at most $O(n^{d-1})$ tuples of total weight $n$.
Its size is not the research claim.

We use the nested convention
\begin{equation}\label{mixed:Li}
 \Li_{s_1,\ldots,s_r}(z_1,\ldots,z_r)
  =\sum_{n_1>\cdots>n_r\ge1}
            \frac{z_1^{n_1}\cdots z_r^{n_r}}
                 {n_1^{s_1}\cdots n_r^{s_r}}.
\end{equation}
The \emph{cumulative colors} are $\xi_j=z_1\cdots z_j$.
The relative colors $z_j$ need not have modulus at most one.
If $\max_j|\xi_j|<1$, the series is locally normally convergent in
all its orders: writing $n_j-n_{j+1}$ as positive increments, with
$n_{r+1}=0$, bounds the color product by $\rho^{n_1}$ for some
$\rho<1$. Polynomial growth from the denominators is then harmless.

\begin{theorem}[Finite mixed-order normal form]\label{mixed:main}
Let $P=\#\{j:a_j>0\}$ and
\[
                \mathcal A=\bigcup_{j=1}^d\{\alpha:\alpha^{p_j}=X_j\}.
\]
The function \eqref{mixed:Z} is a finite linear combination of terms
\begin{equation}\label{mixed:output}
          \Li_{C+v,k_2,\ldots,k_r}(z_1,\ldots,z_r),
 \qquad v\in\Z,\quad k_2,\ldots,k_r\in\Z_{>0},
\end{equation}
whose cumulative colors all belong to $\mathcal A$. The coefficients
are obtained by finite rational operations in the colors and roots,
rational Bernoulli numbers, and integral shuffle multiplicities.
One may choose
\begin{equation}\label{mixed:depth}
 r\le\begin{cases}
 d,&P=d,\\
 P+1,&0<P<d,\\
 1,&P=0.
 \end{cases}
\end{equation}
If all $X_j$ are roots of unity of orders dividing $q_j$, every
cumulative color has order dividing $\lcm_j(p_jq_j)$.

The identity initially holds in the region of absolute convergence
and holds meromorphically for every $C$. If every $X_j\ne1$, both
$Z$ and every final term are entire in $C$. Every fixed derivative
in $C$ is obtained by taking the corresponding derivative with respect
to the first index
of the multiple polylogarithms in \eqref{mixed:output}.
The fixed integral orders $a_j$ are not free differentiation variables
in this formula.
\end{theorem}

The depth in \eqref{mixed:depth} is a bound for this construction;
it is not a minimal-depth or period-independence assertion. In particular,
one nonpositive factor can substantially shorten a high-depth input
when only a few of the remaining indices are positive.

\subsection{Step 1: remove weights and shuffle the positive factors}

Let the coefficient transform be
\[
 \Mcal_C\left[\sum_{n\ge0}f_nt^n\right]
             =\sum_{n\ge1}\frac{f_n}{n^C}.
\]
The constant term is omitted. The generator of \eqref{mixed:Z} is
\begin{equation}\label{mixed:generator}
 K(t)=\prod_j\Li_{a_j}(X_jt^{p_j}),\qquad Z(C)=\Mcal_C K.
\end{equation}
Root filtering gives, for every integer $a$,
\begin{equation}\label{mixed:roots}
            \Li_a(Xt^p)=p^{a-1}\sum_{\alpha^p=X}\Li_a(\alpha t).
\end{equation}
Indeed the sum of $n$th powers of the roots is zero unless $p$ divides
$n$, and is then $pX^{n/p}$. This also checks the factor $p^{a-1}$.
Expanding the product reduces the proof to unit weights and colors
in $\mathcal A$.

For a nonzero letter $\alpha$, use the one-forms
$\omega_0=dt/t$ and $\omega_\alpha=\alpha\,dt/(1-\alpha t)$,
and define iterated integrals recursively by integrating the leftmost
letter against the remaining word. The word
$0^{a-1}\alpha$ represents $\Li_a(\alpha t)$ for $a>0$.
Partitioning the product of the integration simplices by the order
of their coordinates proves the shuffle product, with multiplicities.
Thus the $P$ positive factors give finitely many words with exactly
$P$ nonzero letters. A word with zero-block indices $k_1,\ldots,k_P$
and successive nonzero letters $\beta_1,\ldots,\beta_P$ represents
\begin{equation}\label{mixed:word}
 \Li_{k_1,\ldots,k_P}
       (\beta_1t,\beta_2/\beta_1,\ldots,\beta_P/\beta_{P-1}).
\end{equation}
All the $k_i$ are positive. This identification follows either by
expanding the forms geometrically in the recursive integrals, or by
differentiating and fixing the value at $t=0$. Its cumulative colors
are exactly the word letters, all in $\mathcal A$.

The nonpositive factors form a rational function $R(t)$, since
\[
 \Li_{-M}(z)=(z\,d/dz)^M\frac{z}{1-z}\qquad(M\ge0).
\]
Each factor is bounded at infinity, so there is no positive-degree
polynomial part in the partial-fraction decomposition
\begin{equation}\label{mixed:R}
 R(t)=c_\infty+\sum_{\alpha\in\mathcal A}\sum_{h=1}^{M_\alpha}
                         \frac{c_{\alpha,h}}{(1-\alpha t)^h}.
\end{equation}
An upper bound for $M_\alpha$ is the sum of $1-a_j$ over the
nonpositive factors with color $\alpha$. Even if a numerator cancels
some poles, the valid coefficient formula remains
\begin{equation}\label{mixed:partialcoeff}
 c_{\alpha,h}=[u^{M_\alpha-h}]
                 u^{M_\alpha}R((1-u)/\alpha).
\end{equation}
This finite Taylor extraction handles repeated colors directly.

\subsection{Step 2: a rational prefactor with one free spectral order}

For a nested polylogarithm $H$ as in \eqref{mixed:word}, write its
relative colors as $\beta_1,\ldots,\beta_r$ for this subsection.
Clearly
\[
 \Mcal_C\Li_{k_1,\ldots,k_r}(\beta_1t,\beta_2,\ldots,\beta_r)
   =\Li_{C+k_1,k_2,\ldots,k_r}(\beta_1,\ldots,\beta_r).
\]
Define the rational numbers $b_{h,k}$ by
\begin{equation}\label{mixed:b}
 \binom{v+h-1}{h-1}=
       \prod_{\nu=1}^{h-1}(1+v/\nu)
       =\sum_{k=0}^{h-1}b_{h,k}v^k.
\end{equation}
Their connection to generalized harmonic numbers is explicit:
\[
 \prod_{\nu=1}^{h-1}(1+v/\nu)
   =\exp\left(\sum_{j\ge1}
          \frac{(-1)^{j+1}H_{h-1}^{(j)}}{j}v^j\right).
\]

\begin{lemma}[Rational-prefactor transform]\label{mixed:prefactor}
For $h\ge1$ and $H=\Li_{k_1,\ldots,k_r}
(\beta_1t,\beta_2,\ldots,\beta_r)$,
\begin{align}\label{mixed:prefactorformula}
 \Mcal_C\bigl[(1-\alpha t)^{-h}H\bigr]
 ={}&\Li_{C+k_1,k_2,\ldots,k_r}(\beta_1,\ldots,\beta_r)\notag\\
 &+\sum_{k=0}^{h-1}\sum_{j=0}^k
 b_{h,k}(-1)^j\binom kj\,
 \Li_{C-k+j,k_1-j,k_2,\ldots,k_r}
       (\alpha,\beta_1/\alpha,\beta_2,\ldots,\beta_r).
\end{align}
The first inner index $k_1-j$ is permitted to be nonpositive at this
intermediate step.
\end{lemma}
\begin{proof}
Convolution of coefficients gives a sum over $n\ge m\ge1$ with
factor $\alpha^{n-m}\binom{n-m+h-1}{h-1}$. The diagonal $n=m$
is the first term because the binomial coefficient equals one there.
For $n>m$, expand \eqref{mixed:b} and
$(n-m)^k=\sum_{j=0}^k(-1)^j\binom kj n^{k-j}m^j$.
The new outer color is $\alpha$ and the next is $\beta_1/\alpha$,
which proves the formula by direct reindexing.
\end{proof}

If $P=0$, no $H$ is present. The same binomial coefficient gives
directly
\begin{equation}\label{mixed:rationalonly}
 Z(C)=\sum_{\alpha,h}\sum_{k=0}^{h-1}
                    c_{\alpha,h}b_{h,k}\Li_{C-k}(\alpha),
\end{equation}
the depth-one theorem already proved in the source. Thus the new
construction agrees exactly with that existing sector.

\subsection{Step 3: eliminate every nonpositive inner slot}

\begin{lemma}[Finite elimination, including confluent colors]
\label{mixed:eliminate}
Suppose an inner index in \eqref{mixed:Li}, at position $j\ge2$,
equals $-M$ with $M\ge0$. Put $U=n_{j-1}$ and $L=n_{j+1}$;
for the final position set $L=0$. If $q=z_j\ne1$, define
\[
 R_h(q)=(q\,d/dq)^h\frac1{1-q}.
\]
Then
\begin{equation}\label{mixed:geometric}
 \sum_{n=L+1}^{U-1}q^n n^M
 =\sum_{h=0}^M\binom Mh R_h(q)
       \left\{q^{L+1}(L+1)^{M-h}-q^U U^{M-h}\right\}.
\end{equation}
At $q=1$ the exact replacement is
\begin{equation}\label{mixed:Bernoulli}
 \sum_{n=L+1}^{U-1}n^M
       =\frac{B_{M+1}(U)-B_{M+1}(L+1)}{M+1}.
\end{equation}
In either case the resulting multiple polylogarithms have one less
index. The cumulative alphabet does not enlarge.
\end{lemma}
\begin{proof}
Apply $(q\,d/dq)^M$ to the finite geometric sum
$(q^{L+1}-q^U)/(1-q)$ and use the product rule. This proves
\eqref{mixed:geometric}; the Bernoulli difference identity proves
\eqref{mixed:Bernoulli}.

Expand the powers of $L+1$ into powers of $L$. An upper-end term
multiplies the preceding color by $q$ and subtracts a nonnegative
integer from its order; a lower-end term does the same to the following
color and order. At the last position the lower-end term is a constant.
In every case the $j$th summation disappears. In cumulative colors,
the upper-end operation deletes the old $(j-1)$st cumulative letter,
and the lower-end operation deletes the old $j$th letter. When $q=1$
those two letters coincide. Hence no new cumulative color is introduced.
\end{proof}

If a neighboring positive index becomes nonpositive, repeat the rule.
Depth decreases at every repetition, so the algorithm terminates.
The first index can become any $C+$integer and is never eliminated.
This explains both the precise output class and the depth bound.

\begin{proof}[Completion of Theorem~\ref{mixed:main}]
Apply \eqref{mixed:roots}, the finite shuffle expansion, partial
fractions, Lemma~\ref{mixed:prefactor}, and
Lemma~\ref{mixed:eliminate}. The positive-only case has depth $d$;
a rational prefactor initially adds at most one index to depth $P$;
all subsequent eliminations decrease depth. The rational-only case is
\eqref{mixed:rationalonly}. Alphabet preservation and the root-order
statement follow at each stage. The proof so far uses locally normally
convergent series for interior cumulative colors, and absolute
convergence at unit colors in a sufficiently far right half-plane.

For analytic continuation use
\begin{equation}\label{mixed:Mellin}
 Z(C)=\frac1{\Gamma(C)}\int_0^\infty
                   t^{C-1}\prod_j\Li_{a_j}(X_je^{-p_jt})\dd t.
\end{equation}
It follows by termwise Laplace integration in a common convergent
half-plane. If no $X_j=1$, its kernel is analytic at zero and
exponentially decreasing at infinity. Subtracting an arbitrarily long
Taylor polynomial near zero gives meromorphic terms $1/(C+n)$,
whose poles are canceled by $1/\Gamma(C)$. This proves entire
continuation. The same argument applies to every output term: choose
an integral positive first index as its base Mellin kernel, and note
that its iterated-integral word has no singular letter at $t=1$.
Every cumulative letter lies in $\mathcal A$, which contains no $1$
when no original $X_j=1$. Its first-order continuation is therefore
entire. If some original $X_j=1$, the convergent power--logarithm
expansions below continue the product kernel meromorphically. Each
nested output term has the same type of continuation: a positive-index
iterated word has a convergent local power--logarithm expansion at
$t=1$. This follows recursively from its defining integrals, since
the only singular endpoint form is $\dd t/(1-t)$ and integrating a
convergent power series times a finite power of $\log(1-t)$ preserves
that class. Apply the same Mellin subtraction to a positive first-index
base word to continue its first order. The identity theorem extends
the finite reduction, and differentiation in the
free variable $C$ is legitimate.
\end{proof}

Radial continuation in the coefficient variable means $K(t)\mapsto
K(\rho t)$; in the original variables this is
$X_j\mapsto\rho^{p_j}X_j$, not a common damping of all original
colors. At resonant colors the resulting Abel limit need not be an
ordinary finite value outside the convergence domain. All formulas
there mean the stated meromorphic continuation, with Laurent
coefficients taken afterward.

\subsection{Local Laurent coefficients at every nonpositive outer order}

Let $P_1=\#\{j:a_j>0,\ X_j=1\}$. The integer limits of
Jonqui\`ere's expansion \cite[Eq. 25.12.12]{DLMFPolylog} give
\begin{align}\label{mixed:positiveexpansion}
 \Li_k(e^{-pt})
 ={}&\frac{(-pt)^{k-1}}{(k-1)!}
          \bigl(H_{k-1}-\log p-\log t\bigr)\notag\\
 &+\sum_{n\ge0,\ n\ne k-1}\zeta(k-n)\frac{(-pt)^n}{n!},
       \qquad k\ge1,\\
 \Li_{-M}(e^{-pt})
 ={}&\frac{M!}{(pt)^{M+1}}
          +\sum_{n\ge0}\zeta(-M-n)\frac{(-pt)^n}{n!},
       \qquad M\ge0.\label{mixed:negativeexpansion}
\end{align}
Here $H_0=0$. For $X\ne1$ the ordinary Taylor expansion is
\begin{equation}\label{mixed:colorexpansion}
 \Li_a(Xe^{-pt})=
       \sum_{n\ge0}\frac{(-p)^n}{n!}\Li_{a-n}(X)t^n.
\end{equation}
Products yield a convergent local power--logarithm expansion
\begin{equation}\label{mixed:local}
 \prod_j\Li_{a_j}(X_je^{-p_jt})
       =\sum_{n\ge n_0}\sum_{r=0}^{P_1}
                                  \kappa_{n,r}t^n(\log t)^r.
\end{equation}
At any prescribed $n$, its coefficients require only finitely many
terms of \eqref{mixed:positiveexpansion}--\eqref{mixed:colorexpansion}.

\begin{theorem}[Pole structure and finite local Laurent constants]
\label{mixed:localtheorem}
Set $g(C)=1/\Gamma(C)$. Near $C=-n$, the complete polar contribution
to $Z(C)$ is
\begin{equation}\label{mixed:localmodel}
        g(C)\sum_{r=0}^{P_1}
              \frac{(-1)^r r!\kappa_{n,r}}{(C+n)^{r+1}}.
\end{equation}
The difference from \eqref{mixed:localmodel} is $g(C)$ times a
holomorphic function. For $n\ge0$, in particular,
\begin{equation}\label{mixed:finitepart}
 \boxed{\displaystyle
 \CT_{C=-n} Z(C)=
     \sum_{r=0}^{P_1}\frac{(-1)^r r!}{(r+1)!}
                           \kappa_{n,r}g^{(r+1)}(-n).}
\end{equation}
Possible poles at nonpositive integers have order at most $P_1$;
possible poles at positive integers have order at most $P_1+1$.
There are no other possible poles. Missing local coefficients mean zero.
\end{theorem}
\begin{proof}
Taylor--logarithm subtraction in \eqref{mixed:Mellin} reduces the
local terms to
\[
 \int_0^1t^{C+n-1}(\log t)^r\dd t
                  =\frac{(-1)^r r!}{(C+n)^{r+1}}.
\]
Arbitrarily high subtraction makes the remainder holomorphic on any
prescribed compact set. At $C=-n$ with $n\ge0$, $g$ has a simple
zero. Thus its product with the holomorphic remainder has zero
constant coefficient, and the constant of each displayed local term
is exactly \eqref{mixed:finitepart}. The pole assertions follow from
the same zero and from the fact that the powers $n$ are integers.
\end{proof}

Formula \eqref{mixed:finitepart} does not determine all higher
regular coefficients. For example, the constant coefficient of
$Z'(C)$ generally receives the first coefficient of $g(C)$ times
the global holomorphic remainder. This is the same obstruction that
prevents differentiating an isolated negative-integer value identity
transversely. The finite mixed-order formula, which keeps $C$ free,
does determine such derivatives in its stated multiple-polylogarithm
class.

Reciprocal-Gamma derivatives in \eqref{mixed:finitepart} are finite
polynomials in $\gamma$, zeta values, and harmonic numbers. Indeed
\begin{equation}\label{mixed:reciprocalgamma}
 \frac1{\Gamma(-n+u)}=
       (-1)^n n!\,u\prod_{j=1}^n(1-u/j)
       \exp\left(\gamma u+
             \sum_{k\ge2}\frac{(-1)^{k+1}\zeta(k)}k u^k\right).
\end{equation}
Only degree $r+1$ is needed for a local logarithm of degree $r$.

\subsection{Explicit triple identities and a resonant Laurent expansion}

Define
\begin{equation}\label{mixed:F}
 F(C;z)=\sum_{m,n,k\ge1}\frac{m\,z^{m+n+k}}{nk(m+n+k)^C}.
\end{equation}
For $|z|<1$ it is entire in $C$; at $z=1$ the following original
series converges absolutely for $\Re C>2$. Its finite reduction is
\begin{align}\label{mixed:Fidentity}
 F(C;z)={}&2\Li_{C-1,1,1}(z,1,1)
          -2\Li_{C-1,1}(z,1)+2\Li_{C,1}(z,1)\notag\\
         &+2\Li_{C-1}(z)-2\Li_C(z).
\end{align}
An independent elementary proof groups by $N=m+n+k$:
\[
 \sum_{m+n+k=N}\frac m{nk}
 =2N H_{N-1}(1,1)-2(N-1)H_{N-1}+2(N-1).
\]
Here $H_N(1,1)=\sum_{N\ge j>k\ge1}(jk)^{-1}$.
The formula follows by first putting $v=n+k$, summing
$1/(n(v-n))=2H_{v-1}/v$, and using
$\sum_{v=1}^{N-1}H_{v-1}=(N-1)(H_{N-1}-1)$.
Thus this example tests the general construction independently of
the shuffle implementation.

The coefficient generator is
$\Li_{-1}(z)\Li_1(z)^2$. Consequently
\begin{align}\label{mixed:Fvalues}
 F(0;-1)&=-\frac{\log^2 2}{4},&
 F(-1;-1)&=-\frac{\log2}{4},\notag\\
 F(-2;-1)&=\frac{\log^2 2-\log2-1}{8},&
 F(1;-1)&=\frac{\log^2 2}{2}+\log2-1.
\end{align}
The first three follow by applying $(z\,d/dz)^{-C}$ to the rational
times logarithm generator at the nonpositive integers. For the last,
direct integration gives the stronger functional identity
\begin{equation}\label{mixed:Fprimitive}
 F(1;z)=\int_0^z\frac{\Li_1(u)^2}{(1-u)^2}\dd u
       =\frac{\Li_1(z)^2-2\Li_1(z)+2}{1-z}-2.
\end{equation}
More generally $z\,\partial_z F(C;z)=F(C-1;z)$ and the primitive
of $F(C;z)/z$ vanishing at zero is $F(C+1;z)$. Every fixed
$C$-derivative obeys the same identities.

At $z=1$, the kernel in \eqref{mixed:Mellin} is
\[
 \Li_{-1}(e^{-t})\Li_1(e^{-t})^2
 =\frac{\log^2t}{t^2}-\frac{\log t}{t}
    -\frac{\log^2t}{12}+\frac{\log t}{12}+\frac14
       +O\bigl(t(1+|\log t|^2)\bigr).
\]
Using \eqref{mixed:localmodel} and \eqref{mixed:reciprocalgamma}
therefore proves the exact Laurent expansion
\begin{equation}\label{mixed:Flaurent}
 \boxed{\displaystyle
 F(C;1)=-\frac1{6C^2}-\frac{2\gamma+1}{12C}
       +\frac14-\frac{\gamma^2+\gamma}{12}
                +\frac{\zeta(2)}{12}+O(C).}
\end{equation}
This example has a genuine double pole at a nonpositive outer order.
The finite part is a short Gamma--zeta polynomial, while the whole
function is the depth-three expression \eqref{mixed:Fidentity}.

A second example exercises two different nonpositive-slot eliminations:
\begin{align}\label{mixed:squareexample}
 \sum_{m,n,k\ge1}\frac{k^2z^{m+n+k}}{mn(m+n+k)^C}
 ={}&2\Li_{C-2,1,1}(z,1,1)
       -3\Li_{C-2,1}(z,1)+3\Li_{C-1,1}(z,1)\notag\\
 &+\frac72\Li_{C-2}(z)-\frac92\Li_{C-1}(z)+\Li_C(z).
\end{align}
Its generator is $\Li_{-2}(z)\Li_1(z)^2$. At $z=1$ its finite
value at $C=0$ is $-1/12$. At $C=-1+u$ its Laurent expansion is
\begin{equation}\label{mixed:squareLaurent}
 \frac1{60u^2}+\frac{\gamma/60-13/720}{u}
   +\frac{\gamma^2}{120}-\frac{13\gamma}{720}
    -\frac{\pi^2}{720}-\frac1{480}+O(u).
\end{equation}
For verification, the $t^1$ coefficient of the local kernel is
$-\log^2t/120-\log t/720+1/288$; insertion into
\eqref{mixed:localmodel} at $C=-1$ gives
\eqref{mixed:squareLaurent} directly.

\subsection{A convergent evaluation and exact pole classification}

Write $U(C;z)$ for the left side of \eqref{mixed:squareexample}.
At $z=1$ its defining series converges absolutely for $\Re C>3$.
Grouping by $N=m+n+k$ gives the additional finite harmonic identity
\begin{align}\label{mixed:Ucoeff}
 \sum_{m+n+k=N}\frac{k^2}{mn}
 ={}&N^2\bigl(H_{N-1}^2-H_{N-1}^{(2)}\bigr)
        -3N(N-1)H_{N-1}\notag\\
 &+\frac{7N^2-9N+2}{2}.
\end{align}
This follows by taking the outer coefficient in
\eqref{mixed:squareexample}, including the vanishing cases $N=1,2$.
The following ordinary value uses only classical height-one multiple
zeta identities, whose beta-integral proof is included to fix the
normalization.

\begin{corollary}[A convergent zeta evaluation]
\label{mxex:weight-four}
\begin{equation}
 \boxed{\displaystyle
 \sum_{m,n,k\geq1}\frac{k^2}{mn(m+n+k)^4}
   =\frac{\pi^4}{24}+\frac{7\pi^2}{12}
      -\frac{15}{2}\zeta(3).}
 \label{mxex:convergent-value}
\end{equation}
\end{corollary}

\begin{proof}
Set $C=4,z=1$ in \eqref{mixed:squareexample}, and use
\[
 \zeta(2,1,1)=\zeta(4),\qquad
 \zeta(2,1)=\zeta(3),\qquad
 \zeta(3,1)=\frac14\zeta(4).
\]
For a self-contained verification of these three classical relations,
consider the classical Aomoto--Drinfel'd height-one generating identity,
recorded in Kaneko and Ohno~\cite[Eq.~(5)]{mxex:KanekoOhno2010}:
\begin{equation}\label{mxex:height-one}
 \sum_{r,s\geq1}\zeta(r+1,\{1\}^{s-1})x^ry^s
   =1-\frac{\Gamma(1-x)\Gamma(1-y)}{\Gamma(1-x-y)}.
\end{equation}
Indeed the shuffle identity gives
$(-\log(1-t))^s/s!=\operatorname{Li}_{1,\ldots,1}(t,1,\ldots,1)$,
with $s$ indices.  Expand this absolutely convergent power series and
use
$\int_0^1t^{n-1}(-\log t)^{r-1}\,dt=(r-1)!/n^r$.
Termwise integration is justified by nonnegative summands and gives
\[
 \zeta(r+1,\{1\}^{s-1})
 =\frac1{(r-1)!s!}\int_0^1
   \frac{(-\log t)^{r-1}(-\log(1-t))^s}{t}\,dt.
\]
Summing on a sufficiently small bidisc turns the right side into
\[
 x\int_0^1t^{-x-1}\bigl((1-t)^{-y}-1\bigr)\,dt
 =1+x\frac{\Gamma(-x)\Gamma(1-y)}{\Gamma(1-x-y)}.
\]
The beta evaluation is first justified for $\Re x<0$ and then by
analytic continuation near zero.  The functional equation of Gamma
gives \eqref{mxex:height-one}.  Finally
\[
 \log\frac{\Gamma(1-x)\Gamma(1-y)}{\Gamma(1-x-y)}
 =\sum_{j\geq2}\frac{\zeta(j)}j
   \bigl(x^j+y^j-(x+y)^j\bigr).
\]
Its $xy^2$, $xy^3$, and $x^2y^2$ coefficients give the three displayed
identities; in the last case use $\zeta(2)^2=5\zeta(4)/2$.
Formula \eqref{mixed:squareexample} now gives
$15\zeta(4)/4+7\zeta(2)/2-15\zeta(3)/2$, which is
\eqref{mxex:convergent-value}.
\end{proof}

\subsection{The full pole pattern at the resonant color}

Write $U(C)=U(C;1)$ and let $B_j$ denote the Bernoulli numbers with
$B_1=-1/2$.  For $\Re C>3$,
\begin{equation}\label{mxex:mellin}
 U(C)=\frac1{\Gamma(C)}\int_0^\infty t^{C-1}K(t)\,dt,
 \qquad
 K(t)=\operatorname{Li}_1(e^{-t})^2
                  \operatorname{Li}_{-2}(e^{-t}).
\end{equation}
The integral follows directly from the Gamma integral for the denominator;
absolute convergence follows from the coefficient bound $O(N^2(1+\log N)^2)$.

\begin{proposition}[Actual spectral poles]
\label{mxex:poles}
The meromorphic function $U$ has a triple pole at $C=3$, double poles
at $C=2,1$ and at every negative odd integer, and simple poles at
every negative even integer.  These are all its poles.  The point
$C=0$ is regular and
\begin{equation}\label{mxex:zero-value}
 U(0)=-\frac1{12}.
\end{equation}
For every integer $m\geq1$,
\begin{equation}\label{mxex:even-residue}
 \operatorname*{Res}_{C=-2m}U(C)
       =\frac{m}{m+1}B_{2m+2}.
\end{equation}
For every integer $m\geq0$, the coefficient of the leading double
pole at $C=-(2m+1)$ is
\begin{equation}\label{mxex:odd-leading}
 [(C+2m+1)^{-2}]\,U(C)
       =-\frac{B_{2m+4}}{m+2}.
\end{equation}
\end{proposition}

\begin{proof}
The convergent local expansions for $0<t<2\pi$ are
\begin{align*}
 \operatorname{Li}_1(e^{-t})&=-\log t+a(t),&
 a(t)&=\frac t2-\sum_{r\geq1}
             \frac{B_{2r}}{2r(2r)!}t^{2r},\\
 \operatorname{Li}_{-2}(e^{-t})&=b(t),&
 b(t)&=\frac2{t^3}+\sum_{r\geq1}
             \frac{B_{2r+2}}{(2r+2)(2r-1)!}t^{2r-1}.
\end{align*}
The second expansion follows by differentiating
$(e^t-1)^{-1}$ twice.  The first follows by integrating its negative
and fixing the constant from
$\operatorname{Li}_1(e^{-t})+\log t\to0$.
Thus
\[
 K(t)=b(t)(\log t)^2-2a(t)b(t)\log t+a(t)^2b(t).
\]
Mellin subtraction shows that a coefficient
$\kappa_{n,r}t^n(\log t)^r$ contributes the singular model
\[
 \frac1{\Gamma(C)}
 \frac{(-1)^r r!\kappa_{n,r}}{(C+n)^{r+1}}.
\]
All other contributions are holomorphic locally, after a sufficiently
long Taylor--logarithm subtraction.  There are no additional poles
because the tail of the integral decays exponentially.

At the first four powers,
\begin{equation}\label{mxex:local-first}
 K(t)=\frac{2\log^2t}{t^3}
       -\frac{2\log t}{t^2}
       +\frac{\log t/6+1/2}{t}
       -\frac1{12}+O(t\log^2t).
\end{equation}
The three positive poles therefore have orders three, two, and two,
with nonzero leading coefficients $2$, $2$, and $-1/6$, respectively.
The zero of $1/\Gamma(C)$ at $C=0$ cancels the simple local Mellin
pole there; its derivative at zero is one, proving
\eqref{mxex:zero-value}.

For a positive even $n$, $b(t)$ has no $t^n$ term.  Since the only
odd power in $a(t)$ is $t/2$, the coefficient of $t^n\log t$ in
$K$ is $-b_{n-1}$, where
\[
 b_{n-1}=\frac{B_{n+2}}{(n+2)(n-1)!}.
\]
Using $(1/\Gamma)'(-n)=(-1)^n n!$ gives residue
$n!b_{n-1}=nB_{n+2}/(n+2)$, which is
\eqref{mxex:even-residue}.  For positive odd $n$, the coefficient
of $t^n\log^2t$ is
$b_n=B_{n+3}/((n+3)n!)$.  The same reciprocal-Gamma zero gives
leading double-pole coefficient $-2n!b_n$, proving
\eqref{mxex:odd-leading}.  The even-index Bernoulli numbers appearing
here are nonzero, so these are actual poles, of exactly the stated
orders.
\end{proof}


At the next nonpositive even order one also obtains
\begin{equation}\label{mixed:Uminus2}
 \CT_{C=-2}U(C)=\frac{19}{720}-\frac{\gamma}{60}.
\end{equation}
Indeed $[t^2]K(t)=\log(t)/120+1/1440$, and
$1/\Gamma(-2+u)=2u+(2\gamma-3)u^2+O(u^3)$.
Multiplication by its local Mellin model gives
$2/1440-(2\gamma-3)/120$, proving the displayed constant.
This is a Laurent constant of the continuation, not a value of the
positive-term defining series at $C=-2$.

\subsection{The depth-one boundary: Gamma, polygamma, and Stieltjes data}

When the normal form reaches depth one at a $q$th root of unity
$\alpha$, finite residue-class summation gives
\begin{equation}\label{mixed:DFT}
 \Li_s(\alpha)=q^{-s}\sum_{j=1}^q\alpha^j\zeta(s,j/q).
\end{equation}
For $s\ne1$ its $k$th order derivative is
\begin{equation}\label{mixed:DFTjets}
 q^{-s}\sum_{j=1}^q\alpha^j
  \sum_{\ell=0}^k\binom k\ell(-\log q)^{k-\ell}
                         \zeta^{(\ell)}(s,j/q).
\end{equation}
For nontrivial $\alpha$ the common pole at $s=1$ cancels first,
so the finite formula there is
\begin{equation}\label{mixed:DFTStieltjes}
 \Li_1^{[k]}(\alpha)=\frac1q\sum_{j=1}^q\alpha^j
  \sum_{\ell=0}^k\binom k\ell(-\log q)^{k-\ell}
                         (-1)^\ell\gamma_\ell(j/q).
\end{equation}
These classical identities, already used in the source, fix the
normalization of the depth-one output. Lerch's identity
$\zeta'(0,a)=\log\Gamma(a)-\tfrac12\log(2\pi)$ and
$\psi^{(r)}(a)=(-1)^{r+1}r!\zeta(r+1,a)$, $r\ge1$, give the corresponding
Gamma and polygamma expressions \cite{DLMFHurwitz}. The new theorem
does not erase the positive-depth order derivatives left by its
normal form; further reductions require additional identities.

% END sections/04_mixed.tex


% BEGIN sections/05_sampling.tex
\section{Sampled identities and finite single-centre reduction}
\label{sampling:section}

The mixed-spectral report proves a two-centre functional criterion and
explicitly asks whether identities on integer Fourier modes imply the
corresponding scalar identities \cite[Section 6, question on sampled
identities]{RepoMixed}. In the finite power--logarithm class used there,
the answer is yes. We give an elementary proof that also works on many
sparser sets, then obtain the exact finite-reduction criterion for any
number of distinct twists.

Classical uniqueness theorems, such as Carlson's theorem, concern much
larger holomorphic function classes \cite{Pila}. The elementary lemma
below uses the convergent expansions specific to the present functions.
The contribution to the repository is the complete sampling implication
and its application to finite twisted-family reduction; no new general
Carlson theorem is claimed.

\begin{definition}[Finite power--logarithm germs at infinity]
\label{sampling:class}
Let $\mathscr P$ consist of functions, for sufficiently large positive
$x$, of the form
\begin{equation}\label{sampling:germ}
 f(x)=\sum_{\nu=1}^J x^{\alpha_\nu}
           \sum_{k=0}^{K_\nu}(\log x)^k h_{\nu k}(1/x),
\end{equation}
where the $h_{\nu k}$ are holomorphic in a disc about zero and
$\alpha_\nu\in\C$. Powers of positive $x$ use the real logarithm.
Every finite product or sum of shifted powers and nonnegative integral
powers of their logarithms belongs to this class.
\end{definition}

For the last assertion, expand
$(x+a)^s=x^s(1+a/x)^s$ and
$\Log(x+a)=\log x+\Log(1+a/x)$ for $x>2|a|$.
The expansions are convergent, rather than merely asymptotic.

\begin{theorem}[Uniqueness on an asymptotically dense logarithmic mesh]
\label{sampling:unique}
Let $f\in\mathscr P$. Suppose that $x_j\to\infty$ is strictly
increasing, $x_{j+1}/x_j\to1$, and $f(x_j)=0$ for every sufficiently
large $j$. Then $f$ vanishes identically for all sufficiently large $x$.
Consequently any two finite power--logarithm expressions that agree at
every sufficiently large integer agree as analytic germs at infinity.
They agree on every domain reached by coherent analytic continuation.
\end{theorem}
\begin{proof}
First combine exponents whose differences are integers. In each such
class choose the largest integer translate among the finitely many
displayed exponents; this keeps all coefficient functions holomorphic
at zero. Unless $f$ is identically zero, in each remaining class there
is a first nonzero Taylor coefficient among its finitely many
$h_{\nu k}$. Among these finitely many leading terms take the largest
real exponent $\sigma$, and among terms with that real exponent the
largest logarithmic degree $d$. Convergent Taylor expansion then gives
\begin{equation}\label{sampling:leading}
 \frac{f(x)}{x^\sigma(\log x)^d}
   =\sum_{\ell=1}^q c_\ell e^{\ii\tau_\ell\log x}+o(1),
 \qquad x\to+\infty,
\end{equation}
where the real $\tau_\ell$ are distinct and at least one $c_\ell$ is
nonzero. The remainder tends to zero on the whole positive ray.

Put $P(t)=\sum c_\ell e^{\ii\tau_\ell t}$ and $t_j=\log x_j$.
The hypothesis gives $P(t_j)\to0$ and $t_{j+1}-t_j\to0$.
The derivative of $P$ is bounded; hence interpolation between adjacent
$t_j$ gives $P(t)\to0$ as $t\to\infty$. On the other hand direct
integration of the finitely many exponentials shows
\[
 \lim_{T\to\infty}\frac1T\int_0^T|P(t)|^2\dd t
       =\sum_{\ell=1}^q|c_\ell|^2>0.
\]
This contradicts $P(t)\to0$. Thus the assumed nonzero leading term
cannot occur, and the germ is zero. The final assertion follows from
the identity theorem.
\end{proof}

The logarithmic-mesh condition includes $x_j=j^q$ for every real $q>0$.
It cannot simply be omitted: the nonzero function
\[
 \sin\left(\frac{2\pi\log x}{\log 2}\right)
 =\frac{x^{2\pi\ii/\log2}-x^{-2\pi\ii/\log2}}{2\ii}
\]
belongs to $\mathscr P$ and vanishes at $x=2^j$. Conversely, if
one leaves the finite class entirely, $\sin(\pi x)$ vanishes at every
integer. These examples state exactly what is used in the theorem.

\begin{theorem}[Any number of centres: complete finite reduction]
\label{sampling:centres}
Let $a_1,\ldots,a_d$ be distinct complex numbers and let
$\alpha_1,\ldots,\alpha_d\in\C$. Use coherent branches of
\[
                       f(z)=\prod_{j=1}^d(z+a_j)^{\alpha_j}.
\]
The following are equivalent:
\begin{enumerate}[label=(\roman*)]
\item On a nonempty domain, $f$ is a finite constant-coefficient
linear combination of single-centre order jets
$(z+a_j)^r\Log^k(z+a_j)$, with arbitrary $r\in\C$ and
$k\in\Z_{\ge0}$.
\item Such a finite combination agrees with $f(n)$ at every
sufficiently large positive integer $n$.
\item Either every $\alpha_j$ is integral, or exactly one
$\alpha_j$ is nonintegral and all the other exponents are
nonnegative integers.
\end{enumerate}
The equivalence remains valid if the sampled set in (ii) is any
sequence in Theorem~\ref{sampling:unique}.
\end{theorem}
\begin{proof}
The sampling theorem proves (ii)$\Rightarrow$(i), and analytic
continuation proves the converse after the branches are fixed. If all
exponents are integers, ordinary rational partial fractions give (i).
If only $\alpha_i$ is nonintegral and all remaining exponents are
nonnegative integers, expand their polynomial product in $z+a_i$.
This proves (iii)$\Rightarrow$(i).

For necessity, suppose $\alpha_i\notin\Z$, write $w=z+a_i$, and put
\[
             h(w)=\prod_{j\ne i}(w+a_j-a_i)^{\alpha_j}.
\]
This function is holomorphic near $w=0$. The monodromy difference
$M_i-I$ around a small loop about zero annihilates every single-centre
term based at $-a_j$ with
$j\ne i$. On the product it gives
$(e^{2\pi\ii\alpha_i}-1)w^{\alpha_i}h(w)$.
The remaining finite sum of $w^r\Log^k w$ is annihilated by a
nonzero polynomial $P(E)$ in $E=w\,d/dw$: include a factor
$(E-r)^{k+1}$ for each term. Expanding $h(w)=\sum_{m\ge0}h_mw^m$
therefore yields
\[
                         P(\alpha_i+m)h_m=0\quad(m\ge0).
\]
A polynomial has finitely many roots, so $h$ is a polynomial.
Near every distinct point $w=a_i-a_j$, its factor there is a
nonzero holomorphic factor times $(w+a_j-a_i)^{\alpha_j}$.
Polynomiality forces $\alpha_j\in\Z_{\ge0}$ for every $j\ne i$.
There can be no second nonintegral exponent. This proves necessity.
\end{proof}

\begin{corollary}[Finite closure of several twisted Fourier families]
\label{sampling:twisted}
For distinct real noninteger twists $\theta_j$, let a periodic
distribution $K$ have Fourier coefficients
\[
 \widehat K(n)=\prod_{j=1}^d\bigl(2\pi\ii(n+\theta_j)\bigr)^{\alpha_j},
\]
where $\Log(2\pi\ii x)=\log(2\pi|x|)+\pi\ii\sgn(x)/2$.
Then $K$ is a finite linear combination of single-twist families
and their spectral order derivatives if and only if condition (iii)
of Theorem~\ref{sampling:centres} holds. Allowing an additional finite
Fourier polynomial does not change the necessity of that condition.
\end{corollary}
\begin{proof}
On the positive Fourier tail, every phase and power of $2\pi$ is
a constant. Expanding powers of
$\log(2\pi(n+\theta_j))+\pi\ii/2$ gives precisely the class in
Theorem~\ref{sampling:centres}; a finite Fourier polynomial changes
only finitely many modes. This proves necessity. For sufficiency,
rational partial fractions and polynomial expansion hold coefficient
by coefficient on both tails with the stated branches; only one
factor can be nonintegral. The coefficients have locally uniform
polynomial growth, so Fourier equality proves distributional equality.
\end{proof}

Thus the two-centre criterion in the source extends to the actual
integer Fourier spectrum, and to arbitrarily many centres. The claim
concerns a prescribed finite family of functions. It says nothing
about rational or algebraic independence of individual constants
obtained by evaluating the distributions at particular arguments.

% END sections/05_sampling.tex


% BEGIN sections/06_audit.tex
\section{Source corrections, verification, and exact claim status}
\label{audit:section}

\subsection{Corrections and qualifications for integration}

\paragraph{The independent-order interpolant needs its established attribution.}
Any integration of Section~\ref{eht:sec} should cite Ihara, Nakamura,
and Yamamoto for the one-endpoint interpolant, its entire spectral
continuation, and its harmonic product law \cite{eht:INY}.
Equation~\eqref{eht:INY} identifies the functions exactly. This is a
literature reconciliation that supplies part of an open repository
question, not a newly discovered one-endpoint theorem. The smallest
Gamma-separating kernel requested in the full question remains open.
The Gamma formulas in Section~\ref{eht:sec} evaluate symmetric word
combinations; they do not evaluate each ordering separately.

\paragraph{Scalar-only qualifications can now be strengthened.}
The mixed-spectral source carefully restricts its finite two-centre
criterion to a scalar functional identity and asks about sampled
Fourier coefficients \cite{RepoMixed}. Theorem~\ref{sampling:unique}
supplies exactly the missing implication in its finite power--logarithm
class. Thus that qualification may now be replaced by
Theorem~\ref{sampling:centres} and Corollary~\ref{sampling:twisted}.
This corrects an incomplete research status, rather than a false
existing theorem. No conclusion about individual period independence
should be appended.

\paragraph{Local polylogarithm estimates in a specifically inspected preprint.}
In arXiv:2510.10093v1, equation (2.3) prints a maximum where the local
growth estimate requires a minimum, and equation (2.7) includes $s=1$
in an $O(1)$ case \cite{tr3:SS}. The elementary counterexamples are
\[
 \Li_0(e^{-y})=\frac1{e^y-1}\sim y^{-1},\qquad
 \Li_1(e^{-y})=-\log(1-e^{-y})=-\log y+O(y).
\]
For fixed $r\notin\Z_{>0}$ the corrected first estimate is
\[
 \Li_r(e^{-y})=O\bigl(y^{\min(0,\Re r-1)}\bigr),\qquad y\downarrow0.
\]
For the second, use $O(1+|\log y|)$ at $s=1$; the bounded case
holds for $\Re s\ge1$, $s\ne1$. These statements follow directly
from the convergent local formulas \eqref{mixed:positiveexpansion}
and Jonqui\`ere's expansion.
The subsequent ``any $\epsilon>0$'' in equation (2.8) should be a
sufficiently small positive exponent. For example, its hypotheses
permit any
\[
 0<\epsilon<\Re t+\min(0,\Re s-1),
\]
which also absorbs the logarithm when $s=1$.
These changes repair the local convergence justification. They do not
by themselves refute the paper's continuation identities. The finding
is about this inspected preprint version; the journal version has not
been asserted to contain the same printing errors.

\paragraph{Already reported corrections should not be relabeled as new.}
The directional-zeta package already records the author-PDF corrections
$H_0=0$, the positive-integer logarithm coefficient in the
Borwein--Dilcher local expansion, and the local-series restriction
on the splitting point \cite{RepoDirectional}. The present article
uses the corrected integer Jonqui\`ere formula. It carries those
corrections forward by reference instead of claiming a second discovery.

\paragraph{Spectral restriction, finite parts, and color confluence.}
The Tornheim origin is not a joint regular point. Keep the polar
quotients in \eqref{tr3:model} until a direction has been specified.
The cubic theorem excludes $a+c=0$ and $b+c=0$; it must not be
evaluated on either polar hyperplane by cancelling a numerical zero.
The mixed-order theorem keeps the outer order $C$ free, so its
$C$-derivatives are valid. It keeps the inner integers fixed, so their
transverse derivatives are not supplied. At a resonant color take
the meromorphic germ before its Laurent coefficients. Apparent
denominators in a nonconfluent partial-fraction expression are handled
by the repeated-pole Taylor formula or by the Bernoulli branch of the
elimination rule, not by substituting $q=1$ in \eqref{mixed:geometric}.

\subsection{The Gaussian conjectures and the bounded unsuccessful search}

At the pinned revision, the relevant canonical labels are
\src{cycloquot:conj:S6} and \src{s8new:conj:S8} \cite{RepoMain}.
The earlier rejected $S_8$ candidate is a different statement; it
must not be confused with the revised candidate. The already proved
$S_4$ identity retains its proved status. Nothing in this article
proves or refutes the current $S_6$ or $S_8$ candidates.

A bounded exploratory $S_6$ calculation was attempted with the frozen
primitive weight-seven target. It generated 51,244 distinct imaginary
double-shuffle rows for factor-weight splits $(1,6)$, $(2,5)$, and
$(3,4)$, including the common single-divergence cancellation, and two
convergent Cayley rows. The system uses 24,514 imaginary word
coordinates. Elimination modulo $2,147,483,647$ was interrupted before
the relation system was exhausted: the last checkpoint had 8,500
pivots, 29,853,917 stored entries, and 393 remaining target terms.

Those numbers report an unfinished computation, not a rank theorem.
No row-membership certificate, no nonmembership result, and no
separating functional for that enlarged system was obtained. The
generator, letter conventions, target, and termination record are
included under \src{exploratory/s6}. Large matrices and compiled
executables are omitted because the generator reproduces them.
A future successful modular search would still require rational
reconstruction and independent exact replay before proving the
conjecture. A future negative conclusion would have to specify an
exhausted relation space and certify its separating functional.

\subsection{Reproducible verification and what it establishes}

The archive contains a single runner, \src{code/verify_all.py}, and
separate scripts for the three computational tracks. Exact checks
use integers, rational numbers, or exact algebraic expressions.
The mixed-order checker evaluates original weighted coefficients by
direct convolution and compares them with finite nested sums from
the output; it does not reuse the shuffle or elimination code to
construct its reference values.

\begin{table}[htbp]
\centering\small
\begin{tabular}{@{}L{0.23\textwidth}L{0.46\textwidth}L{0.24\textwidth}@{}}
\toprule
Track & Independent checks & Recorded outcome \\
\midrule
Hurwitz interval functions
& Finite strict sums, ordered quotient, reversed-star form, residue
  samples, Bernoulli reduction, and a corrupted-sign control
& 743 exact assertions; 340 words included \\
\addlinespace
Hurwitz analytic diagnostics
& Euler--Maclaurin continuation, singular-intersection Cauchy
  extraction, shifts, composition, Gamma products, and quadrature
& 24 diagnostics at 65 digits; largest residual about
  $3.03\cdot10^{-31}$ \\
\addlinespace
Cubic Tornheim directions
& Symbolic quadratic constraints and cancellation identities;
  independent subtracted Mellin evaluations on three rays
& Exact checks pass; largest numerical discrepancy about
  $2.27\cdot10^{-38}$ at 65 digits \\
\addlinespace
Mixed-order compiler
& Original coefficient convolution against output nested sums,
  including repeated colors, depth four, and weighted roots
& 192 exact coefficients in 13 varied cases \\
\addlinespace
Explicit triple specializations
& Independent elementary kernel expansions, Laurent coefficients,
  Bernoulli leading poles, Euler primitive, and convergent Mellin integral
& 24 exact assertions; $C=4$ residual about $3.89\cdot10^{-61}$
  at 60 digits \\
\bottomrule
\end{tabular}
\caption{Verification scope. Decimal residuals are diagnostics, not
interval-certified accuracy claims.}\label{audit:checks}
\end{table}

The mathematical proofs establish arbitrary depth, arbitrary allowed
directions, and meromorphic continuation. No finite collection of
coefficient checks establishes those generalities by itself. The
compiler is a transparent implementation of the finite proof;
its algebraic simplifier is not claimed to have optimal complexity.
Inputs with floating-point colors or uncertified color domains are
rejected rather than silently treated as exact numbers.

The package also records the source commit, hashes of the incoming
archives and consulted source files, the individual result ledgers,
and the PDF build and visual inspection. The modular TeX sections
have distinct label prefixes \src{eht:}, \src{tr3:}, \src{mixed:},
and \src{sampling:}. Integration notes describe dependencies and
conjecture statuses. The standalone TeX is generated from those
modular sources so that the PDF can be rebuilt without the ProveIt
repository or a network connection.

% END sections/06_audit.tex


% BEGIN sections/07_research.tex
\section{Further exact research questions}
\label{research:section}

The questions below concern identities, finite reductions, and the
structure of the remaining analytic coordinates. They are proposed
problems, not conjectures asserted on numerical evidence. Each specifies
an initial case or a concrete output that could settle the problem.

\begin{question}[Evaluate or relate the two cubic coordinates]
Find a second functional identity that reduces $\Omega$ or $\kappa$
to a prescribed algebra of ordinary zeta derivatives, Gamma or Barnes
values, and elementary constants. The convergent moment in
\eqref{tr3:kappa} and the polylogarithmic integral
\eqref{tr3:Dpolylog} provide two explicit starting points. An equality
between those representations is not itself an ordinary-constant
evaluation. A useful first target is a reflection or multiplication
identity for the whole Dirichlet series $D(A)$, whose second derivative
controls $\kappa$.
\end{question}

\begin{question}[The quartic and higher directional layers]
The cubic theorem uses the degree-two Taylor jet of the holomorphic
remainder $R$. At degree four, compute the degree-three jet subject
to symmetry, the exact $C$ axis, and the Euler relation
\eqref{tr3:stuffleR}. Give an explicit finite set of convergent Gamma
or multiple-Gamma moments spanning all $W^{(4)}_{a,b,c}(0)$.
Determine the exact dimension of the formal jet quotient under the
proved functional relations. This is a finite linear-algebra question
and should be separated from numerical or arithmetic independence
of the resulting constants.
\end{question}

\begin{question}[Systematic cancellation of directional coordinates]
The coefficient of $\kappa$ vanishes on a conic, and cyclic averaging
eliminates it identically. At each higher derivative order, construct
the polynomial map from directions to the residual jet coordinates.
Compute its kernel for permutation sums and for small finite families
of rational directions. The desired output is an exact criterion
for combinations reducible to already known coordinates, with the
smallest number of directions needed within that specified formal
relation space.
\end{question}

\begin{question}[The minimal Gamma-separated independent-order kernel]
Complete Question 8 of the exact-identities report: beyond the entire
interpolant identified here, is there a smallest explicit generating
kernel whose singular part separates through a Gamma factor when
all inner orders move independently? Begin at depth three near
$(s_1,s_2,s_3)=(1,1,1)$ and retain three independent perturbations.
Use the coefficients $C_{k,m}$ in \eqref{eht:C} to compare candidate
factorizations at every intersecting prefix divisor. State the
equivalence relation under which ``smallest'' is meant.
\end{question}

\begin{question}[Nonsymmetric harmonic jets and the coordinate $\eta$]
Equations \eqref{eht:first-eta}--\eqref{eht:second-eta} show precisely
where $\eta(a)$ survives after entire normalization. Determine which
linear combinations of second and third independent order jets cancel
all new ordered coordinates. At rational endpoints, compute a finite
character-projection matrix including conductor logarithms. A first
target is the antisymmetric second jet at depth two with endpoints
$a=1/2$, $b=1$; an ordinary Stieltjes reduction would require an
additional identity beyond the symmetric Gamma generator.
\end{question}

\begin{question}[Reflection and multiplication of normalized primitives]
For the primitive \eqref{eht:primitive} and its nonpositive-inner
extensions, derive exact formulas under $a\mapsto1-a$ and
$a\mapsto qa$. Keep the base endpoint fixed or explicitly transform
it, so additive constants remain determined. Differentiate the whole
spectral identity to obtain log-weighted Gamma and Stieltjes
primitives. Compare the resulting constants with standard Barnes
multiple-Gamma normalizations at the first two derivative orders.
\end{question}

\begin{question}[Minimal depth within the mixed-order output class]
The bound $P+1$ is constructive, not optimal in every example.
Characterize when a mixed input reduces to depth at most $P$, or to
ordinary Hurwitz functions and their order derivatives. Start with
one negative factor and two positive factors at a fixed root-of-unity
alphabet. Any minimality theorem must specify the coefficient field,
the allowed argument transformations, and the set of functional
relations; otherwise a depth count alone is not an obstruction.
\end{question}

\begin{question}[A finite map for character and color cancellations]
Apply root-of-unity projections to the exact compiler before evaluating
periods. Can the positive-depth terms be represented in a finite
quotient whose kernel detects all reductions to the depth-one
Stieltjes formulas \eqref{mixed:DFTStieltjes} at a fixed weight and
derivative order? Repeated colors should be treated within the
confluent formula, so the reduction matrix remains meaningful when
several input colors coincide. An exact row certificate is the
appropriate output.
\end{question}

\begin{question}[Spectral derivatives beyond locally determined constants]
Theorem~\ref{mixed:localtheorem} makes every Laurent constant at a
nonpositive outer integer a local calculation. The next regular
coefficient usually depends on the global Mellin remainder. Classify
linear combinations of mixed kernels for which that remainder cancels
to one or two prescribed orders. The two explicit triple kernels
in Section~\ref{mixed:section} provide a small initial family;
their finite MPL normal forms permit legitimate differentiation
before specialization.
\end{question}

\begin{question}[Simultaneous color and order resonance]
Let several $X_j$ tend to $1$ while $C$ tends to a pole and its
derivatives are taken. Determine a joint completed germ whose local
color logarithms and spectral Laurent coefficients are explicit.
The radial variable must respect the weights,
$X_j\mapsto\rho^{p_j}X_j$. A first target is one negative and two
positive inner orders with two colors approaching $1$ at different
rates. Distinguish a joint analytic continuation from an ordinary
Abel limit of a divergent series.
\end{question}

\begin{question}[Four-factor kernels and higher coincident products]
Use the mixed-order compiler for the integer-index pieces of the
four-factor Fourier zero-frequency constraint. Identify a finite
collection of constrained zeta kernels for the different sign
partitions, then compute the fully coincident fourth digamma finite
part in one fixed endpoint coordinate. The cubic directional formulas
should supply boundary checks. The unresolved global kernel must not
be replaced by an undefined product of singular distributions.
\end{question}

\begin{question}[Sampling sets beyond vanishing logarithmic gaps]
Theorem~\ref{sampling:unique} is sharp against unrestricted geometric
sampling, but a particular set of imaginary exponents may admit much
sparser uniqueness sets. For a prescribed finite frequency set
$\{\Im\alpha_\nu\}$, characterize the sampling sequences on which
its trigonometric leading part is determined. Separate this finite
frequency problem from extensions to asymptotic expansions with
exponentially small terms, where functions invisible to every
power--logarithm coefficient can occur.
\end{question}

\begin{question}[Exact Gaussian certificates]
Resume $S_6$ from the frozen rational target and the precisely specified
relation space, using a strategy that controls sparse fill and records
row provenance. A modular candidate must be lifted and replayed over
the rationals. If the chosen space does not contain the target,
compute a separating functional for that exhausted space and then
identify a mathematically justified new relation family. Treat the
revised $S_8$ candidate separately; numerical agreement alone cannot
promote either to a theorem.
\end{question}

\begin{question}[Formal proof kernels and reusable interfaces]
Formalize the finite root filter, sparse shuffle, repeated-pole Taylor
extraction, geometric/Bernoulli elimination, and polynomial jet maps.
Their analytic interface consists of a reusable Mellin-subtraction
lemma with locally normal parameter dependence and an explicit
reciprocal-Gamma zero. A second interface should formalize prefix
residue cancellation at intersecting divisors. These finite and
analytic components can then support later repository identities
without conflating exact algebra, analytic continuation, and
numerical evaluation.
\end{question}

% END sections/07_research.tex


% BEGIN references.tex
\phantomsection
\addcontentsline{toc}{section}{References}
\begin{thebibliography}{99}
\small
\raggedright

\bibitem{RepoMain}
Vladimir Reshetnikov, \emph{ProveIt: Polylogarithms manuscript},
repository revision \texttt{5a790187c8e186e41e2b990b4941cb7a1a3c7b6b},
inspected 11 October 2026.
\url{https://github.com/VladimirReshetnikov/ProveIt/tree/5a790187c8e186e41e2b990b4941cb7a1a3c7b6b/Analysis/Polylogarithms/docs/manuscript}.
Canonical conjecture labels \src{cycloquot:conj:S6} and
\src{s8new:conj:S8}.

\bibitem{RepoIncoming}
ProveIt, \emph{Incoming research reports}, eleven archives at the
same pinned revision.
\url{https://github.com/VladimirReshetnikov/ProveIt/tree/5a790187c8e186e41e2b990b4941cb7a1a3c7b6b/docs/incoming}.
The accompanying provenance ledger lists individual filenames and hashes.

\bibitem{RepoExact}
ProveIt incoming report, \emph{Exact Identities},
\src{ProveIt_Exact_Identities_2026-10-10.zip}, at the pinned revision.
Especially \src{sections/audit_research.tex}, Question 8,
``Independent inner exponents.'' Available in \cite{RepoIncoming}.

\bibitem{RepoMixed}
ProveIt incoming report, \emph{Mixed Spectral Identities},
\src{ProveIt_Mixed_Spectral_Identities_2026-10-10.zip}, at the pinned
revision. Especially the harmonic and unequal-twist sections and
\src{sections/06_audit_research.tex}, Question 7,
``From scalar functional reduction to sampled identities.''
Available in \cite{RepoIncoming}.

\bibitem{RepoOrdered}
ProveIt incoming report, \emph{Ordered Hurwitz Resonances},
\src{ProveIt_Ordered_Hurwitz_Resonances_2026-10-10.zip}, at the pinned
revision. Especially the common-shift ordered Laurent coefficient
$\eta(a)$ and its shift and differentiation formulas.
Available in \cite{RepoIncoming}.

\bibitem{RepoDirectional}
ProveIt incoming report, \emph{Coincident Stieltjes and Directional Zeta},
\src{ProveIt_Coincident_Stieltjes_Directional_Zeta_2026-10-10.zip},
at the pinned revision. Especially Sections 2, 5--7, the external-source
audit, and Section 9, Questions 2, 8, 11, and 13.
Available in \cite{RepoIncoming}.


% BEGIN references_hurwitz.tex
\bibitem{eht:INY}
K. Ihara, Y. Nakamura, and S. Yamamoto,
\emph{Interpolant of truncated multiple zeta functions},
The Ramanujan Journal \textbf{67} (2025), article 1.
\url{https://doi.org/10.1007/s11139-025-01045-2}.
Preprint: \url{https://arxiv.org/abs/2407.20509}.
Especially Proposition 2.1, Corollary 4.2, and Section 5.

\bibitem{eht:MehtaViswanadham}
J. Mehta and G. K. Viswanadham,
\emph{Analytic continuation of multiple Hurwitz zeta functions},
Journal of the Mathematical Society of Japan \textbf{69} (2017), 1431--1442.
\url{https://doi.org/10.2969/jmsj/06941431}.

\bibitem{eht:Hoffman}
M. E. Hoffman,
\emph{Quasi-shuffle products},
Journal of Algebraic Combinatorics \textbf{11} (2000), 49--68.
\url{https://doi.org/10.1023/A:1008791603281}.
Preprint: \url{https://arxiv.org/abs/math/9907173}.

% END references_hurwitz.tex


% BEGIN references_tornheim.tex
% Bibliography entries for cubic_rays.tex; include inside thebibliography.
\bibitem{tr3:incoming}
ProveIt incoming report,
\emph{Coincident Stieltjes and Directional Zeta},
\texttt{ProveIt\_Coincident\_Stieltjes\_Directional\_Zeta\_2026-10-10.zip},
Sections 2, 6, and 9,
repository revision \texttt{5a790187c8e186e41e2b990b4941cb7a1a3c7b6b}.
\url{https://github.com/VladimirReshetnikov/ProveIt/tree/5a790187c8e186e41e2b990b4941cb7a1a3c7b6b/docs/incoming}.

\bibitem{tr3:BD}
Jonathan M.~Borwein and Karl Dilcher,
\emph{Derivatives and fast evaluation of the Tornheim zeta function},
The Ramanujan Journal \textbf{45} (2018), 413--432.
\href{https://doi.org/10.1007/s11139-017-9890-9}{doi:10.1007/s11139-017-9890-9}.
Inspected author version, titled
\emph{Derivatives and fast evaluation of the Witten zeta function}:
\url{https://carmamaths.org/resources/jon/omega3.pdf}.

\bibitem{tr3:DLMF}
NIST Digital Library of Mathematical Functions,
\S\S5.9, 5.11, 25.11, and 25.12.
\url{https://dlmf.nist.gov/}.

\bibitem{tr3:SS}
Sumukha Sathyanarayana and N.~Guru Sharan,
\emph{Mordell--Tornheim zeta function: Kronecker limit type formulas
and special values}, Advances in Applied Mathematics
\textbf{177} (2026), article 103075.
\href{https://doi.org/10.1016/j.aam.2026.103075}{doi:10.1016/j.aam.2026.103075}.
Inspected preprint: arXiv:2510.10093v1 (2025).
\url{https://arxiv.org/html/2510.10093v1}.

\bibitem{tr3:Dobrowolski}
Przemys\l aw Dobrowolski,
\emph{Evaluation of the symmetrized Mordell--Tornheim zeta function},
arXiv:2603.20550v1 (2026).
\url{https://arxiv.org/html/2603.20550v1}.

% END references_tornheim.tex


% BEGIN references_examples.tex
% Include these entries within the article's thebibliography environment.
\bibitem{mxex:KanekoOhno2010}
M.~Kaneko and Y.~Ohno,
\emph{On a kind of duality of multiple zeta-star values},
International Journal of Number Theory \textbf{6} (2010), no.~8,
1927--1932, doi:10.1142/S179304211000385X.
Author manuscript: \url{https://www2.math.kyushu-u.ac.jp/~mkaneko/papers/45zeta-star.pdf}.
Equation~(5) records the height-one Gamma generating identity and attributes
it to Aomoto and Drinfel'd.

% END references_examples.tex


\bibitem{DLMFPolylog}
NIST Digital Library of Mathematical Functions, \S25.12,
\emph{Polylogarithms}. Especially equations 25.12.11 and 25.12.12.
\url{https://dlmf.nist.gov/25.12}.

\bibitem{DLMFHurwitz}
NIST Digital Library of Mathematical Functions, \S25.11,
\emph{Hurwitz Zeta Function}. Especially the order-zero derivative,
parameter derivative, and special-value formulas.
\url{https://dlmf.nist.gov/25.11}.

\bibitem{Panzer}
Erik Panzer, \emph{Algorithms for the symbolic integration of
hyperlogarithms with applications to Feynman integrals},
Computer Physics Communications \textbf{188} (2015), 148--166.
\url{https://arxiv.org/abs/1403.3385}.
\url{https://doi.org/10.1016/j.cpc.2014.10.019}.

\bibitem{Pila}
Jonathan Pila, \emph{Note on Carlson's theorem},
Rocky Mountain Journal of Mathematics \textbf{35} (2005), no. 6,
2107--2112.
\url{https://research-information.bris.ac.uk/en/publications/note-on-carlsons-theorem/}.

\end{thebibliography}

% END references.tex

\end{document}
